% =============================================================================
% MANUSCRIPT: Discrimination from Soft Signals: Dilemmas of Mimic Men
% =============================================================================
% Author: Anurag Srivastava (Riskcare Ltd., London; University of Reading)
% ORCiD:  0000-0002-6477-4430
%
% FIRST DRAFT — generated from handover.tex, proposal.tex, and working session
% =============================================================================

\documentclass[12pt,a4paper]{article}

%% Packages
\usepackage[a4paper, margin=1.1in]{geometry}
\usepackage{amsmath, amssymb, amsthm}
\usepackage{setspace}
\usepackage{booktabs}
\usepackage{array}
\usepackage{hyperref}
\usepackage{natbib}
\usepackage{graphicx}
\usepackage{enumitem}
\usepackage{titlesec}
\usepackage{parskip}
\usepackage{microtype}
\usepackage{xcolor}

%% Spacing and layout
\onehalfspacing
\setlength{\parindent}{0pt}
\setlength{\parskip}{8pt}

%% Theorem environments
\newtheorem{proposition}{Proposition}
\newtheorem{theorem}{Theorem}
\newtheorem{lemma}{Lemma}
\newtheorem{corollary}{Corollary}
\newtheorem{definition}{Definition}
\newtheorem{remark}{Remark}
\theoremstyle{definition}
\newtheorem{example}{Example}

%% Title formatting
\titleformat{\section}{\large\bfseries}{\thesection.}{0.5em}{}
\titleformat{\subsection}{\normalsize\bfseries}{\thesubsection.}{0.5em}{}
\titleformat{\subsubsection}{\normalsize\itshape}{\thesubsubsection.}{0.5em}{}

%% Hyperref
\hypersetup{colorlinks=true, linkcolor=blue!60!black,
            citecolor=blue!60!black, urlcolor=blue!60!black}

%% Custom commands
\newcommand{\E}{\mathbb{E}}
\newcommand{\R}{\mathbb{R}}
\newcommand{\G}{\mathcal{G}}
\newcommand{\M}{\mathcal{M}}
\newcommand{\C}{\mathcal{C}}
\newcommand{\A}{\mathcal{A}}

%% ---------------------------------------------------------------
\begin{document}

\begin{center}
{\LARGE\bfseries Discrimination from Soft Signals:\\[0.4em]
Dilemmas of Mimic Men}\\[1.5em]

{\large Anurag Srivastava}\\[0.3em]
{\normalsize Riskcare Ltd., London $\cdot$ University of Reading}\\[0.2em]
{\normalsize ORCiD: 0000-0002-6477-4430}\\[1em]

{\normalsize \today}\\[1.5em]

\begin{minipage}{0.85\textwidth}
\small\noindent\textbf{Abstract.}
Social boundaries seem to have persisted in spite of individual
counter-examples, legal equality, and formal anti-discrimination
instruments. A prominent economic framework for understanding this
persistence is \citet{Loury2002}, which distinguishes
``discrimination in contract'' (formal, labour-market) from
``discrimination in contact'' (informal, social) and argues that
the latter sustains racial inequality even after the former is
removed. This paper extends Loury's framework to explain how
non-discriminatory social environments amidst profit incentives
could together shape social discrimination in a market of soft
signals.
When people read soft signals---accents, dress, consumption
patterns, phenotype---the updating rule that governs their
inference is better represented with Jeffrey conditioning rather
than standard Bayes. Under this rule, the statistical updating
does not revise a fixed group-level belief: individual
counter-examples are absorbed without changing the group
stereotype. The boundary re-asserts itself through the inferential
structure of soft-signal processing, not through irrationality.
The discrimination mechanism has two sufficient statistics: marker
salience (how visible group markers are) and the group prior (the
sediment of economic history). Because these enter the mechanism
at different points, the instruments that affect each are
structurally different and neither can compensate for the other.
The model generates endogenous closure equilibria that are
path-dependent on economic history and marker salience. Class-based
social closure exists in a racially homogeneous society as a
structural feature; racial discrimination is a perturbation of this
underlying mechanism. There exists a threshold level of
intergenerational mobility---increasing in marker salience---below
which the closure trap persists permanently and above which it is
eliminated through a saddle-node bifurcation. Individual
compliance with host-group markers is not merely futile but
actively harmful through adverse selection. Forced marker
suppression produces reallocation, not elimination. Profit-maximising
platforms calibrate atmosphere using legally invisible instruments,
and it is possible that stricter anti-discrimination law may
constrain formal channels while raising the return to informal
ones. The model is positioned as a formal microfoundation for
Loury's ``discrimination in contact'' programme, with his
framework recovered as a special case corresponding to the
standard-Bayesian limit of the Jeffrey conditioning
mechanism.\\[0.5em]

\noindent\textbf{JEL codes:} D83, J15, J71, D62, Z13.\\
\textbf{Keywords:} Jeffrey conditioning, soft signals,
discrimination in contact, closure equilibrium, marker salience,
intergenerational mobility, boundary persistence.
\end{minipage}
\end{center}


% =============================================================================
\section{Introduction}
\label{sec:intro}
% =============================================================================

Social boundaries persist. They persist despite individual
counter-examples, despite legal equality, despite every formal
instrument aimed at dissolving them. Individual members of
stigmatised groups achieve economic success, adopt host-group
markers, and are personally reclassified---yet the group boundary
re-asserts itself in the next encounter. \citet{Loury2002}
provided a prominent framework for understanding this persistence,
distinguishing ``discrimination in contract'' (formal,
labour-market) from ``discrimination in contact'' (informal,
social) and arguing that the latter sustains inequality even after
the former is removed. But the inferential mechanism that sustains
discrimination in contact---why individual counter-examples fail
to erode group beliefs---was not formally derived. This paper
provides it, showing how non-discriminatory social environments
amidst profit incentives could together shape social discrimination
in a market of soft signals.

When the operative signals in social interaction are soft---accents,
dress, consumption patterns, phenotype---the theoretically
appropriate updating rule is Jeffrey conditioning rather than
standard Bayes. Jeffrey conditioning revises beliefs about group
membership while holding group-level class distributions fixed.
The consequence is a rigidity condition under which individual
counter-examples are absorbed without revising the group prior.
The boundary re-asserts itself as a theorem. The phenomenon is
recognised outside economics: Naipaul's Ralph Singh
\citep{Naipaul1967} invests everything in metropolitan markers and
is personally reclassified while his group is not; Fanon's
colonised subject discovers that the white mask does not revise
the white gaze \citep{Fanon1952}; Coetzee's post-apartheid figures
navigate a world whose legal structure has changed overnight while
its inferential architecture has not \citep{Coetzee1999}. What
these writers identified as a dilemma, the model derives as an
equilibrium outcome.

Three observations motivate the formal model.

The first is economic. Friends of lower economic background cannot
share the same restaurants, leisure activities, schools, or holiday
destinations as friends of higher economic background. This is not
discrimination---it is non-overlapping budget sets producing natural
social segregation. Economic class determines social closure to a
large extent, and it does so in every society, regardless of racial
composition.

The second is inferential. Social discrimination is not only about
economic incompatibility. It is explicitly about race, ethnicity,
religion, language---markers that carry economic connotations through
historical association. One assesses another person's likely economic
position not by observing their bank balance but by reading soft
signals: accent, dress, consumption patterns, name, neighbourhood.
These signals are probabilistically correlated with economic class
through the group's economic history. This is where subjective
probability enters. The question is: what is the correct formal
inference rule for updating beliefs about economic class from soft
social markers?

The third is institutional. Formally non-discriminatory institutions
nevertheless produce discriminatory social compositions. An elite
restaurant imposes no racial bar, yet its clientele is
overwhelmingly from one demographic. The restaurant manages room
composition through informationally mediated instruments: dress
codes calibrated by class, reservation systems filtering by social
network, location creating transport barriers. Legal prohibition of
explicit discrimination appears to make this worse by pushing
discrimination into a shadow market of social information---coded
language in reviews, word of mouth in networks, selective social
media.

We argue that the mechanism connecting these three observations is
\emph{Jeffrey conditioning}---the framework of probability
kinematics developed by \citet{Jeffrey1965}. In Jeffrey conditioning,
an agent revises their probability distribution over a partition
$\{g_1, \ldots, g_K\}$ (group memberships) without receiving a
definitive signal about which partition element is true. The revision
is driven by the evidential weight of a soft marker, not by a
likelihood ratio. Crucially, the conditional probabilities
$F_{g_k}$---the group-level priors over economic class---are
held fixed during the revision. Only the weights $P'(g_k)$ change.

This rigidity condition is the formal expression of the
``exception to the rule'' phenomenon: an individual counter-example
(a high-class member of a stigmatised group) shifts $P'(g_k)$---the
observer attributes the individual to a different group---but does
\emph{not} update $F_{g_k}$ for the group as a whole. The
stereotype persists. This emerges as a theorem of the updating rule,
not as an assumption about preferences or cognitive bias.

The paper makes five contributions. First, we provide a formal
probability-theoretic microfoundation for Loury's
``discrimination in contact'' programme
\citep{Loury2002}, showing that Loury's framework is a special case
corresponding to the standard-Bayesian limit of the Jeffrey mechanism
(Appendix~\ref{app:loury}). The normative urgency of Loury's
programme---that discrimination in contact is a serious injustice
requiring remedy---is not in question here; what is in question is
the formal structure of the mechanism he identifies. Second, we prove that the Jeffrey
updating rule generates a self-sustaining class closure trap as a
stable equilibrium outcome even in a racially homogeneous society
(Proposition~\ref{prop:class-closure}): the trap is path-dependent,
intergenerationally persistent, and resistant to individual
counter-examples by construction. That economic class generates
social clustering is a premise; that it generates a structurally
stable trap is a theorem. Racial discrimination enters as a
perturbation of this already-operating equilibrium, not a separate
phenomenon. Third, we derive a threshold level of intergenerational
mobility $\gamma^*$ above which the closure trap is eliminated
through a saddle-node bifurcation
(Proposition~\ref{prop:bifurcation}), providing a formal
equilibrium-selection justification for mobility programmes. Fourth,
we show that profit-maximising platforms create tokenism as a
structural outcome (Proposition~\ref{prop:tokenism}), that stricter
anti-discrimination law raises the return to informal discrimination
networks (Proposition~\ref{prop:shadow}), and that the welfare loss
decomposes into four components of which one is structurally
irreducible under any policy preserving soft social signals
(equation~\ref{eq:welfare}, Table~\ref{tab:welfare}). Fifth, we
derive an endogenous marker investment result
(Appendix~\ref{app:veblen}) connecting the model to Veblen's theory
of conspicuous consumption and Spence's signalling framework.

The paper is structured as follows. Section~\ref{sec:lit} positions
the model relative to existing theories of discrimination.
Section~\ref{sec:model} develops the base model: primitives, Jeffrey
conditioning, bilateral updating, and comparative statics.
Section~\ref{sec:closure} establishes the closure equilibrium,
intergenerational dynamics, and the mobility threshold.
Section~\ref{sec:platform} develops the platform extension: experience
utility, atmosphere vectors, shadow information markets, platform
exclusion, and tokenism. Section~\ref{sec:discussion} provides
worked examples including the restaurant case, delivery-driver
credential checks, cross-group proximity misreading, compounding
markers in institutional credibility assessment
(Section~\ref{subsec:compounding}), and the
colonial compliance mechanism (Section~\ref{subsec:colonial}).
Section~\ref{sec:empirical} outlines the empirical strategy.
Section~\ref{sec:conclusion} concludes.

% =============================================================================
\section{Related Literature}
\label{sec:lit}
% =============================================================================

The model departs from the existing discrimination literature along
four dimensions: domain, signal structure, updating rule, and prior
formation.

\subsection{Taste-based and statistical discrimination}

\citet{Becker1957} models discrimination as an exogenous taste
parameter---a utility cost of interacting with members of a
particular group. The employer pays for discrimination through
reduced profit. The model predicts that competitive markets erode
discrimination over time, a prediction that has not been borne out
empirically.

\citet{Phelps1972} introduces statistical discrimination: employers
use group membership as a signal about worker productivity when
individual information is costly. \citet{Arrow1973} shows that
statistical discrimination can be self-confirming---if employers
invest less in training workers from a stigmatised group, the group's
productivity falls, confirming the prior. The self-confirming
equilibrium is the closest antecedent to our closure equilibrium,
but Arrow's mechanism operates through the labour market, uses hard
group signals, and relies on standard Bayesian updating.

Our model operates in the domain of social interaction, where there
is no profit motive to discipline incorrect beliefs. It uses soft
signals (accents, dress, consumption) rather than hard group
membership. And it uses Jeffrey conditioning rather than standard
Bayes, which generates the rigidity condition absent from all
existing models.

\subsection{Coate--Loury and affirmative action}

\citet{CoateLoury1993} show that affirmative action can sustain
``patronising equilibria'' in which the policy itself lowers the
incentive for workers from the targeted group to invest in skills.
Our model provides a different---and complementary---justification
for affirmative action: temporary forced integration shifts the
system from the high-discrimination to the low-discrimination
equilibrium, after which the policy is no longer needed
(Proposition~\ref{prop:class-closure}). The policy works by
equilibrium selection, not by compensating for disadvantage.

\subsection{The Akerlof--Spence limit and the hard/soft signal distinction}
\label{subsec:akerlof-spence}

The most natural information-economics response to discrimination
is Spencian signalling: if marginalised individuals credibly signal
their type, the pooling equilibrium breaks down and discrimination
erodes. This logic rests on Bayesian updating over hard signals---a
university degree, a certification, a warranty---where the signal
partitions the type space and the observer's posterior moves with
each observation. Applied to social discrimination, it implies that
marginalised communities should be able to erode group priors by
signalling their true economic quality more assertively. This is the
implicit logic of respectability politics and of the self-confirming
equilibria in \citet{CoateLoury1993}: the signalling equilibrium is
self-fulfilling, so policy intervention that shifts beliefs unlocks
the good equilibrium.

The present model identifies precisely why this does not transfer.
In \citet{Akerlof1970} and \citet{Spence1973}, the signal is hard:
it determines a definite evidential relationship with the type,
and Bayesian updating on hard signals aggregates cleanly from
individuals to the group prior. In social interaction, the operative
signals are soft---accent, dress, consumption patterns---and the
theoretically appropriate rule for soft signals is Jeffrey
conditioning, not Bayesian updating. The rigidity condition
(Proposition~\ref{prop:graded}) decouples individual
reclassification from group prior revision: the compliant individual
is reclassified, the group prior $\bar{c}_g$ does not move.
No volume of individual signalling can drive $\bar{c}_g$ across the
mobility threshold $\bar{c}_g^{\mathrm{mid}}$, because the Jeffrey
architecture holds conditional probabilities fixed by construction.
The Spence fix fails here not because of irrationality, prejudice,
or coordination failure, but because the inference rule governing
soft-signal environments is categorically different from the one
governing hard-signal markets.

\subsection{Loury's programme and the Jeffrey formalisation}

\citet{Loury2002} distinguishes \emph{discrimination in contract}
(formal, now legally constrained) from \emph{discrimination in
contact} (informal social avoidance, legal, consequential for
economic mobility). He identifies biased social inference and racial
stigma as the mechanism sustaining discrimination in contact; his
formal models take this inferential bias as given rather than
deriving it from the updating rule itself.

Our model provides the inferential microfoundation Loury's programme
calls for. The Jeffrey conditioning mechanism is the formalisation of
``biased social inference''; the rigidity condition is the
formalisation of why individual counter-examples fail to update group
priors. In Appendix~\ref{app:loury} we show that Loury's model is
the $\rho \to 0$ limit of our framework, where $\rho$ parameterises
the degree of Jeffrey rigidity. At $\rho = 0$ (standard Bayes),
individual counter-examples fully update the group prior, the
informational welfare loss $\mathcal{L}_{\text{info}}$ vanishes, and
the closure trap disappears---recovering Loury's single-equilibrium
result. The closure trap requires $\rho > \rho^*$, a threshold we
compute explicitly.

Jeffrey conditioning is not merely a modelling choice. It is the
theoretically natural updating rule when signals are soft: it is the
unique rule that revises the marginal distribution over groups while
holding the conditional distributions $F_{g_k}$ fixed ---
precisely the rigidity condition the model requires. Standard
Bayesian updating requires a likelihood function
$P(\text{signal} \mid \text{state})$, which presupposes that the
signal determines a definite evidential relationship with the state.
Soft markers do not have this property. An accent does not determine
group membership; it shifts the evidential weight across groups
without resolving it. Alternative frameworks for updating under
imprecise evidence---Dempster-Shafer belief functions, lower
previsions \citep[see][]{Walley1991}---do not deliver conditional
rigidity as a structural property of the revision rule and therefore
cannot generate the closure trap as an equilibrium outcome.
\citet{Jeffrey1983} calls Jeffrey conditioning ``Bayesianism with a
human face''---an updating rule that respects the genuine uncertainty
of soft evidence.

The model takes the initial group-level priors $\bar{c}_{g_k}$
as historically given. It is silent on whether those priors were
generated by class structures that preceded racial categorisation or
by racial exploitation that created class stratification---a question
that the historical literatures of \citet{DuBois1935},
\citet{BonillaSilva1997}, and racial formation theory address at
length. The formal mechanism operates identically under either
account: whatever the historical process that deposited the initial
priors, Jeffrey conditioning sustains them against individual
counter-examples and generates the same closure trap. The model is
therefore consistent with either historical ordering and takes no
position on which account of the origins of racialised class
structure is correct.

\subsection{Other related work}

\citet{AkerlofKranton2000} introduce identity as a direct argument in
the utility function. Our model does not require identity preferences;
discrimination emerges from inference, not taste. \citet{BordaloEtal2016}
model stereotypes through representativeness heuristics. Our mechanism
is structurally different: Jeffrey conditioning is a coherent
probability revision, not a bias, and the rigidity condition is a
theorem of the updating rule, not a cognitive limitation.

\citet{CurrariniEtal2009} model homophily in networks as the outcome
of strategic choice. Our bilateral updating structure
(Section~\ref{subsec:bilateral}) is complementary: agents in our
model form beliefs about potential network partners through Jeffrey
conditioning, and the resulting engagement decisions generate the
homophilous network structures that Currarini et al.\ analyse.

The Bernoulli prior structure connecting class outcomes within a
group to the De Finetti representation theorem
(Section~\ref{subsec:intergenerational}) places the model in the
subjectivist probability tradition of De~Finetti, Jeffrey, and Savage.
The connection between conspicuous consumption and class signalling
draws on \citet{Veblen1899} and is formalised in
Appendix~\ref{app:veblen} using \citet{Spence1973}'s signalling
framework.

\subsection{The industrial organisation channel}

This paper approaches discrimination from a consumption and
industrial organisation perspective, rather than from the identity
and self-perception channel that dominates the sociological and
behavioural-economics literatures. The question is not how
discrimination affects the self-image, aspiration, or social identity
of those discriminated against --- important as those questions are
--- but how business incentives interact with discrimination as a
market equilibrium phenomenon. Specifically: do profit-maximising
firms reinforce or erode discrimination, under what market conditions,
and through what mechanisms?

The existing IO literature on discrimination is sparse. \citet{Becker1957}
predicts that competitive markets erode taste-based discrimination
because discriminating firms face higher costs. The empirical
record does not support this prediction, and the statistical
discrimination literature explains why: the discrimination can be
individually rational, eliminating the cost disadvantage. Neither
framework asks how firm conduct shapes the \emph{belief architecture}
that sustains discrimination --- the subject of this paper's platform
extension.

\citet{Galbraith1952} argues that market power generates countervailing
power on the other side of the market: as a stigmatised group's
aggregate purchasing power grows, corporate interests partially align
with their inclusion. In the model this corresponds to the exclusion
threshold $\bar{c}_g^{\text{excl}}$ (Proposition~\ref{prop:exclusion}):
as the group's prior rises, the platform's profit motive eventually
generates voluntary partial integration. But countervailing power
deposits the group in the tokenism zone
(Proposition~\ref{prop:tokenism}) --- platform access without social
integration. Market forces alone cannot push the group past the social
tipping point $\bar{c}_g^{\text{mid}}$. This simultaneously vindicates
and limits Galbraith's optimism: the market does generate partial
inclusion, but inclusion falls structurally short of integration.

Relatedly, \citet{Galbraith1967} argues that large corporations
exercise a planning function that shapes rather than merely responds
to demand, using advertising and social conditioning to manage client
preferences. The atmosphere vector $\mathbf{a}$ in
Section~\ref{sec:platform} is the formal counterpart: the platform
does not passively respond to existing client beliefs about social
composition; it actively calibrates those beliefs through the
physical and social environment it creates. Discrimination in this
setting is not a market failure in the standard sense --- it is a
profit-maximising outcome of a firm with market power over its
clients' informational environment.

\begin{table}[t]
\centering
\renewcommand{\arraystretch}{1.3}
\small
\begin{tabular}{lllll}
\toprule
\textbf{Model} & \textbf{Domain} & \textbf{Signal} &
\textbf{Updating} & \textbf{Prior} \\
\midrule
Phelps (1972)           & Labour  & Hard group    & Std Bayes        & Exogenous \\
Arrow (1973)            & Labour  & Hard group    & Std Bayes        & Endogenous (eq.) \\
Akerlof--Kranton (2000) & General & Identity      & None (pref.)     & Exogenous \\
Bordalo et al.\ (2016)  & General & Group stats   & Representativeness & Exogenous \\
Loury (2002)            & Social  & Social markers & Biased (informal) & Partially endog. \\
\textbf{This paper}     & \textbf{Social} & \textbf{Soft markers} &
\textbf{Jeffrey conditioning} & \textbf{Endogenous} \\
\bottomrule
\end{tabular}
\caption{Positioning relative to existing models of discrimination.}
\label{tab:positioning}
\end{table}

% =============================================================================
\section{The Base Model}
\label{sec:model}
% =============================================================================

\subsection{Primitives}

Consider a continuum of agents indexed by $i$ and $j$. When two
agents encounter each other, agent $j$ observes agent $i$'s soft
markers and forms a belief about $i$'s economic class; simultaneously,
agent $i$ observes agent $j$'s soft markers and forms a belief about
$j$'s economic class. Engagement is decided bilaterally on the basis
of these beliefs. The model's engagement decisions are grounded
throughout in economic compatibility: racial markers, gender signals,
cultural indicators, and other soft cues enter the engagement decision
only insofar as they are informative about economic class through the
group prior. This is not a claim that agents consciously reason about
income; it is a claim that the latent variable mediating social
engagement is economic class, and soft markers are its observable
correlates.

Each agent is characterised by:
\begin{itemize}[noitemsep]
  \item Economic class $c_i \in \C = [0,1]$, where $c = 0$ is the
  minimum and $c = 1$ the maximum. Determines consumption bundle and
  is the latent variable underlying the engagement decision.
  \item Social marker $m_i \in \M$, an observable characteristic
  (accent, dress, name, consumption choice). Correlated with $c_i$
  but does not determine it. This is a \emph{soft signal}.
  \item Group membership $g_i \in \G = \{g_1, \ldots, g_K\}$.
  Group membership may be directly observable (race, visible gender)
  or latent and inferred from markers (class, ethnicity, religion).
  In either case what matters for the model is not group membership
  itself but the economic class distribution conditional on group
  membership, which is always a soft inference.
  \item Self-belief $Q_i(c)$, agent $i$'s belief over their own
  socially relevant economic position.
\end{itemize}

Two agents are \textbf{economically compatible} if
$d(c_i, c_j) \leq \bar{d}$, where $d(c_i, c_j) = (c_i - c_j)^2$ is
squared class distance and $\bar{d}$ is the compatibility threshold.
Squared distance is adopted throughout; it decomposes expected
incompatibility into a bias term (distance between group means) and
a variance term (within-group dispersion), a separation that drives
the comparative statics of Section~\ref{subsec:interaction}.

\subsection{The prior and its endogenous formation}
\label{subsec:prior}

Agent $j$'s prior over the economic class of a member of group $g$:
\begin{equation}
  \bar{c}_g = \frac{1}{T} \sum_{t=1}^{T} c_{g,t}
  \label{eq:prior}
\end{equation}
where $c_{g,t} \in [0,1]$ is the class outcome of the $t$-th
historical member of group $g$ and $T$ is the historical window.
The group mean $\bar{c}_g$ is a \textbf{sufficient
statistic of the group's economic history}. Groups with histories of
marginalisation---slavery, colonial extraction, persistent
poverty---have low $\bar{c}_g$.

\begin{remark}[Beta--De~Finetti structure]
\label{rem:bernoulli}
If individual class outcomes within a group are modelled as
exchangeable $[0,1]$-valued random variables, the De~Finetti
representation theorem \citep{DeFinetti1964} guarantees a mixing
measure over possible group means. The natural parametric family
is the Beta distribution $\mathrm{Beta}(\bar{c}_g \nu,\,
(1 - \bar{c}_g)\nu)$ with mean $\bar{c}_g$ and variance
\begin{equation}
  \sigma_g^2 = \frac{\bar{c}_g(1 - \bar{c}_g)}{1 + \nu}
  \label{eq:beta-var}
\end{equation}
where $\nu > 0$ is a concentration parameter governing within-group
heterogeneity. As $\nu \to 0$, the Beta distribution concentrates
on $\{0,1\}$ and the model reduces to the Bernoulli special case
with variance $\bar{c}_g(1-\bar{c}_g)$, which is the parametrisation
adopted throughout the remainder of the paper. This variance plays
a central role in the intergenerational dynamics
(Section~\ref{subsec:intergenerational}).
\end{remark}

\subsection{Jeffrey conditioning: the updating mechanism}
\label{subsec:jeffrey}

Agent $j$ observes soft marker $m_i$. This shifts $j$'s distribution
over group membership from $P(g_k)$ to $P'(g_k)$ without determining
which group $i$ belongs to. Under \textbf{Jeffrey conditioning}:
\begin{equation}
  P'(c) = \sum_{k=1}^{K} F_{g_k} \cdot P'(g_k)
  \label{eq:jeffrey}
\end{equation}

The revision $P \to P'$ is driven by the evidential weight of the
marker, not a likelihood ratio.

\textbf{Rigidity condition.} The conditional probabilities
$F_{g_k}$ are held fixed during the Jeffrey revision. Only
$P'(g_k)$ changes. Consequence: an individual counter-example shifts
$P'(g_k)$---the individual is attributed to a different group---but
does \emph{not} update $F_g$. This is the formal
expression of the ``exception to the rule'' phenomenon and stereotype
persistence. It emerges as a theorem of Jeffrey conditioning, not as
an assumption.

\subsection{Bilateral updating and the reflexive information structure}
\label{subsec:bilateral}

The information structure is bilateral and reflexive. Both agents
update simultaneously. Agent $i$'s self-belief is revised through
observing $j$'s consumption signals:
\begin{equation}
  Q_i'(c_i) = \sum_{k} Q_i(c_i \mid c_j = c_k) \cdot R_i'(c_k)
  \label{eq:bilateral}
\end{equation}
where $R_i'(c_k)$ is $i$'s Jeffrey-revised belief about $j$'s class
after observing $j$'s consumption bundle.

Three consequences follow. First, exclusion suppresses aspirations
and self-perception, not just economic opportunities
(\emph{self-signalling and aspiration formation}). Second, agent $i$
has incentive to manipulate markers to shift $P'(g_k)$ in $j$'s
Jeffrey revision---the social analogue of \citet{Spence1973}'s
signalling model (\emph{endogenous signal production};
Appendix~\ref{app:veblen}). Third, excluded agents have worse
information about available economic opportunities; the closure trap
is also an information trap (\emph{closure as epistemic trap}).

\subsection{The interaction decision}
\label{subsec:interaction}

Social engagement utility:
\begin{equation}
  U(c_i, c_j) = v - \kappa \cdot d(c_i, c_j)
  \label{eq:utility}
\end{equation}
where $v > 0$ is baseline interaction value, $\kappa > 0$ is
incompatibility sensitivity, and $d(\cdot)$ is the class distance.
Outside option: $\bar{u}$.

Expected utility under Jeffrey-updated posterior:
\begin{equation}
  \E[U \mid m_i] = v - \kappa \cdot \E[d(c_i, c_j) \mid m_i]
  \label{eq:exp-utility}
\end{equation}

Agent $j$ engages iff:
\begin{equation}
  \E[d(c_i, c_j) \mid m_i] \leq
  \frac{v - \bar{u}}{\kappa} \equiv \Delta^*(c_j)
  \label{eq:threshold}
\end{equation}

Under bilateral structure, engagement requires condition
(\ref{eq:threshold}) to hold for both agents simultaneously.

\textbf{Two-group simplification.} For the remainder of the
analysis, we work with two groups $\G = \{g_1, g_2\}$ (host and
immigrant/out-group respectively).

\textbf{Marker salience.} Not all markers are equally informative
about group membership. A parameter $s_g \in [0,1]$, which we call
\emph{marker salience}, governs the strength of the Jeffrey
revision---the probability that the observer attributes the
encountered individual to the out-group upon observing their
markers:
\begin{equation}
  P'(g_2 \mid m_i) = s_g
  \label{eq:salience}
\end{equation}
Marker salience is a property of the \emph{marker--observer
interaction}, not an intrinsic property of the group. A South Asian
individual in England has high $s_g$ (phenotypic markers strongly
identify out-group membership); the same individual in South Asia
has $s_g \approx 0$. Mutable markers (accent, dress, consumption)
can be modified, reducing $s_g$ over time; immutable markers (skin
colour, facial structure) produce permanently high $s_g$.

Under this parametrisation, the expected squared distance for an
agent $j$ with class $c_j$ encountering an individual whose markers
trigger salience $s_g$ is:
\begin{equation}
  \E[(c_i - c_j)^2 \mid m_i]
  = s_g \underbrace{\left[(\bar{c}_{g_2} - c_j)^2 + \sigma_{g_2}^2\right]}_{D_2}
  + (1 - s_g) \underbrace{\left[(\bar{c}_{g_1} - c_j)^2 + \sigma_{g_1}^2\right]}_{D_1}
  \label{eq:Ed-salience}
\end{equation}
where $D_k = (\bar{c}_{g_k} - c_j)^2 + \sigma_{g_k}^2$ is the
expected squared distance under certain attribution to group $g_k$.
When $s_g = 1$, the observer attributes the individual to the
out-group with certainty and the expression reduces to:
\begin{equation}
  \E[(c_i - c_j)^2 \mid g_2]
  = (\bar{c}_{g_2} - c_j)^2 + \sigma_{g_2}^2
  \label{eq:Ed-group}
\end{equation}
which is the workhorse expression of the model. When $s_g = 0$,
the markers are invisible and the observer treats the individual
as a member of the host group: no discrimination occurs regardless
of economic history.

The engagement threshold is violated when either the bias (class
distance between observer and group mean) or the uncertainty
(within-group dispersion), weighted by marker salience, is too
large.

\begin{remark}[Space-time sufficiency]
\label{rem:spacetime}
The two sufficient statistics of the discrimination mechanism
correspond to the spatial and temporal dimensions of social
observation. Marker salience $s_g$ is a \emph{spatial} property:
it is determined by what the observer perceives at the point of
encounter---phenotype, dress, accent, body language---and governs
\emph{whether} the Jeffrey revision fires and how strongly.
The group prior $\bar{c}_g$ is a \emph{temporal} property: it is
the sediment of the group's economic history across generations
and governs \emph{what} the Jeffrey revision maps to once it
fires. Every soft signal that enters the model is either a spatial
property of the individual observable at the encounter or encodes
temporal information accumulated across the group's economic
history. The model's central result---that sustained discrimination
requires both channels to be active---can therefore be restated:
sustained discrimination requires both spatial identifiability
($s_g > 0$) and temporal disadvantage ($\bar{c}_g <
\bar{c}_g^{\mathrm{mid}}$). Spatial interventions (reducing
marker visibility through assimilation or ``passing'') and temporal
interventions (raising $\gamma$ above $\gamma^*$ through mobility
programmes) are categorically different and cannot substitute for
each other.
\end{remark}

\subsection{Comparative statics}

\begin{proposition}[Income effect]
\label{prop:income}
Higher own economic class increases the expected class proximity to
other high-class agents, expanding the engagement set:
$\partial \E[(c_i-c_j)^2\mid g]/\partial c_j = -2(\bar{c}_g - c_j) < 0$ when $c_j < \bar{c}_g$.
\end{proposition}

\begin{proposition}[Education effect]
\label{prop:education}
Higher own education increases precision of marker-to-group mapping,
increasing responsiveness of $P'(g_k)$ to individual information and
reducing weight on group prior $F_g$.
\end{proposition}

\begin{proposition}[Economic history effect]
\label{prop:history}
Lower $\bar{c}_g$ generates higher posterior expected
incompatibility for any given marker, pushing more potential partners
below the engagement threshold:
$\partial \E[(c_i-c_j)^2\mid g]/\partial \bar{c}_g = 2(\bar{c}_g - c_j) < 0$ when
$c_j > \bar{c}_g$.
\end{proposition}

\begin{proposition}[Graded discrimination]
\label{prop:graded}
The Jeffrey framework predicts continuous variation in avoidance with
marker ambiguity. A faint accent produces less avoidance than a
strong one. $\E[(c_i-c_j)^2]$ is quadratic in $\bar{c}_g$
($\partial^2 \E/\partial \bar{c}_g^2 = 2 > 0$), distinguishing the model
from Phelps--Arrow's binary response to hard group membership.
\end{proposition}

\begin{proposition}[Marker salience effect]
\label{prop:salience}
Higher marker salience amplifies expected incompatibility when the
out-group is economically disadvantaged:
\[
  \frac{\partial \E[(c_i-c_j)^2 \mid m_i]}{\partial s_g}
  = D_2 - D_1 > 0
  \quad \text{when } \bar{c}_{g_2} < \bar{c}_{g_1}.
\]
Two groups with identical economic histories $\bar{c}_g$ face
different closure dynamics if their marker salience differs. Marker
salience and economic history are jointly necessary for sustained
discrimination: $s_g = 0$ eliminates discrimination regardless of
$\bar{c}_g$, and $\bar{c}_{g_2} = \bar{c}_{g_1}$ eliminates
discrimination regardless of $s_g$.
\end{proposition}

\begin{proposition}[Forced marker suppression]
\label{prop:suppression}
Let $\lambda_s \in [0,1]$ be a policy instrument that caps effective
marker salience: $s_g^{\mathrm{eff}} = \min(s_g,\, \lambda_s)$.
When the cap binds ($\lambda_s < s_g$):
\begin{enumerate}[noitemsep]
  \item The discrimination rate falls:
  $F(\bar{c}_g,\, \lambda_s) < F(\bar{c}_g,\, s_g)$.
  \item The mobility threshold falls:
  $\tilde{\gamma}^*(\lambda_s) < \tilde{\gamma}^*(s_g)$.
  \item The shadow information value of the remaining unsuppressed
  markers rises: $\partial \Omega_j / \partial \lambda_s < 0$
  conditional on $\lambda_s < s_g$. Suppressing one class of markers
  raises the discriminatory return to those that remain.\footnote{The
  model uses a single aggregate $s_g$. The substitution claim is a
  structural argument: if the model were extended to a vector
  $\mathbf{s} = (s_1, \ldots, s_K)$ with channel-specific
  suppression, partial suppression on channel $k$ would raise
  the shadow return to unsuppressed channels. A formal
  multi-channel extension is left for future work.} Partial
  suppression produces marker reallocation, not elimination.
  \item When suppression is lifted ($\lambda_s \to s_g$), the
  equilibrium snaps to the state the unsuppressed system would
  have converged to. If $\bar{c}_g$ has not crossed
  $\bar{c}_g^{\mathrm{mid}}$ during the suppression period, the
  closure trap re-emerges immediately. The speed of
  re-stratification is proportional to the gap
  $s_g - \lambda_s$.
\end{enumerate}
\end{proposition}

% =============================================================================
\section{The Closure Equilibrium}
\label{sec:closure}
% =============================================================================

\subsection{The dynamic system}

If fraction $\phi^t$ of agents discriminate against group $g$ in
generation $t$, the next generation's prior evolves as:
\begin{equation}
  \bar{c}_g^{t+1} =
  \underbrace{\bar{c}_g^t(1 - \delta \phi^t)}_{\text{exclusion erodes prior}}
  + \underbrace{\mu(1 - \bar{c}_g^t)}_{\text{natural convergence}}
  + \underbrace{\tilde{\gamma}\,\bar{c}_g^t(1-\bar{c}_g^t)}_{\text{mobility from within-group dispersion}}
  \label{eq:dynamic}
\end{equation}
where $\delta > 0$ is the economic cost of social exclusion per unit
of discrimination, $\mu \in (0,1)$ is the exogenous intergenerational
convergence rate, and $\gamma \in (0,1)$ is the intergenerational
mobility parameter governing regression to the group mean.

Under the Beta--De~Finetti structure
(Remark~\ref{rem:bernoulli}), the within-group variance is
$\sigma_g^2 = \bar{c}_g(1-\bar{c}_g)/(1+\nu)$, so the mobility term
becomes $\tilde{\gamma}\,\bar{c}_g(1-\bar{c}_g)$ where
$\tilde{\gamma} = \gamma/(1+\nu)$ is the effective mobility
parameter. This term (Remark~\ref{rem:bernoulli}). It is maximised at
$\bar{c}_g = 1/2$ (maximum heterogeneity) and vanishes at $\bar{c}_g = 0$
and $\bar{c}_g = 1$ (homogeneous group). This term captures
intergenerational regression to the group mean: children of
high-class families face some probability of downward mobility, and
children of low-class families have some chance of upward mobility.

The discrimination rate $\phi^t = s_g \cdot F(\bar{c}_g^t)$ is
determined by the engagement threshold (\ref{eq:threshold}),
modulated by marker salience $s_g$
(Proposition~\ref{prop:salience}): when $s_g = 0$, markers are
invisible and no discrimination fires. The composed map is:
\begin{equation}
  H(\bar{c}_g;\, s_g) = \mu + \bar{c}_g(1 - \mu + \tilde{\gamma}
  - s_g\,\delta\, F(\bar{c}_g))
  - \tilde{\gamma}\,\bar{c}_g^2
  \label{eq:composed}
\end{equation}

The boundary conditions $H(0; s_g) = \mu > 0$ and
$H(1; s_g) = 1$ hold for all $s_g$, since $F(1) = 0$ and the
$\tilde{\gamma}$ terms cancel at $\bar{c}_g = 1$.

\begin{definition}[Closure Equilibrium]
\label{def:closure}
A closure equilibrium is a fixed point $(\phi^*, \bar{c}_g^*)$ of the
composed map such that beliefs are consistent with outcomes and
outcomes with beliefs.
\end{definition}

\subsection{Fixed point existence and multiplicity}

\begin{proposition}[Class closure in a homogeneous society]
\label{prop:class-closure}
In the single-group model ($\G = \{g_1\}$), markers signal class
position directly. For sufficiently small $\gamma$ and sufficiently
large $\delta$, the map $H$ has at least two stable fixed points
$\bar{c}_g^L < \bar{c}_g^H = 1$, separated by an unstable tipping point
$\bar{c}_g^{\mathrm{mid}}$. The low fixed point constitutes a
within-group class closure trap: a segment of the population is
permanently excluded from high-class social spaces across
generations, without any inter-group dynamic.
\end{proposition}

\begin{proof}[Proof sketch]
Define $\Psi(\bar{c}_g) = H(\bar{c}_g) - \bar{c}_g$. Then $\Psi(0) = \mu > 0$
and $\Psi(1) = 0$ (since $F(1) = 0$ and the $\gamma$ terms cancel).
By continuity and the Intermediate Value Theorem, $\Psi$ has a root
in $[0,1]$. The multiplicity condition (the S-shape of $H$) ensures
$\Psi$ dips below zero on $(0,1)$, producing two interior roots: a
stable low fixed point and an unstable tipping point. Local stability
follows from $|H'(\bar{c}_g^L)| < 1$ and $|H'(\bar{c}_g^{\mathrm{mid}})| > 1$.
Path dependence follows from the monotonicity of orbits on each side
of $\bar{c}_g^{\mathrm{mid}}$. \qed
\end{proof}

\subsection{The intergenerational mobility threshold}
\label{subsec:intergenerational}

\begin{proposition}[Mobility threshold and bifurcation]
\label{prop:bifurcation}
There exists $\gamma^* > 0$ such that:
\begin{enumerate}[noitemsep]
  \item If $\gamma < \gamma^*$: the closure trap persists (two stable
  fixed points).
  \item If $\gamma > \gamma^*$: the closure trap is eliminated
  (unique stable fixed point at $\bar{c}_g = 1$).
  \item At $\gamma = \gamma^*$: a saddle-node bifurcation---the low
  fixed point and the tipping point collide and annihilate.
\end{enumerate}
\end{proposition}

\begin{proof}[Proof sketch]
At $\gamma = 0$, the closure trap exists
(Proposition~\ref{prop:class-closure}). As $\gamma$ increases, the
mobility term raises $H(\bar{c}_g)$ in the interior, shrinking
the region where $\Psi < 0$. At $\gamma = \gamma^*$, the region
vanishes: the low fixed point and the tipping point merge,
$H'(\bar{c}_g^*) = 1$, and the system undergoes a saddle-node
bifurcation. For $\gamma > \gamma^*$, $\Psi > 0$ on all of $(0,1)$
and the only fixed point is $\bar{c}_g = 1$. The exact $\tilde{\gamma}^*$ is
found numerically by solving the bifurcation system
$\{H(\bar{c}_g) = \bar{c}_g,\; H'(\bar{c}_g) = 1\}$ jointly.
\qed
\end{proof}

The intergenerational persistence of class position near the
bifurcation is characterised in Appendix~\ref{app:Tstar}.

\begin{corollary}[Marker salience raises the mobility threshold]
\label{cor:salience-threshold}
The critical mobility threshold $\tilde{\gamma}^*$ is strictly
increasing in marker salience:
$\partial \tilde{\gamma}^* / \partial s_g > 0$. At $s_g = 0$,
$\tilde{\gamma}^* = 0$ and no closure trap exists for any positive
mobility rate. At $s_g = 1$, $\tilde{\gamma}^*$ attains its maximum,
recovering the base-model threshold. Consequently, groups with
immutable markers (high $s_g$) require strictly higher
intergenerational mobility to escape the closure trap than groups
with mutable markers (low $s_g$), even when starting from identical
economic positions.
\end{corollary}

\begin{proof}[Proof sketch]
The composed map $H(\bar{c}_g;\, s_g)$ is decreasing in $s_g$ on the
interior of $[0,1]$ (since $\partial G/\partial s_g = -\delta\phi
\bar{c}_g < 0$). A higher $s_g$ lowers $H$ in the interior, deepening
the region where $\Psi < 0$ and requiring higher $\tilde{\gamma}$ to
restore $\Psi \geq 0$ everywhere. At $s_g = 0$, the discrimination
channel is fully shut off ($\phi^t = 0$ for all $t$), so
$H(\bar{c}_g;\, 0) = \mu + \bar{c}_g(1 - \mu + \tilde{\gamma}) -
\tilde{\gamma}\bar{c}_g^2$, which satisfies $\Psi \geq 0$ for all
$\tilde{\gamma} \geq 0$. Numerical verification confirms strict
monotonicity across the full range of $s_g$. \qed
\end{proof}

\subsection{The enemy group perturbation}

\begin{definition}[Enemy group]
Group $g_2$ is an enemy group relative to host group $g_1$ if
$\bar{c}_{g_2}^0 < \bar{c}_g^{\mathrm{mid}}$---their prior starts
in the high-discrimination basin.
\end{definition}

The enemy group's integration trajectory depends on two dimensions:
economic history ($\bar{c}_{g_2}^0$) and marker salience ($s_g$).
A group with low $\bar{c}_{g_2}^0$ but also low $s_g$---phenotypically
similar to the host, with mutable cultural markers---may escape
the trap rapidly because the discrimination channel fires weakly.
A group with comparable $\bar{c}_{g_2}^0$ but high $s_g$---visibly
distinct, with immutable phenotypic markers---faces a structurally
harder path. The Irish in England illustrate the former: despite a
history of famine, colonial extraction, and armed conflict,
phenotypic similarity ($s_g \approx 0.15$) meant the Jeffrey
revision fired weakly, and integration occurred within two
generations. Indian and Kenyan immigrants illustrate the latter:
comparable or better initial economic positions but high phenotypic
salience ($s_g \approx 0.9$) sustains the closure trap
indefinitely under the same mobility parameters.

\begin{proposition}[Enemy group trap]
\label{prop:enemy}
If $g_2$ enters with $\bar{c}_{g_2}^0 < \bar{c}_g^{\mathrm{mid}}$
and $s_g$ sufficiently high that $\tilde{\gamma} <
\tilde{\gamma}^*(s_g)$, then $g_2$ converges to the
high-discrimination equilibrium $\bar{c}_g^L(g_2)$ independently
of $g_1$'s trajectory, and $g_1$'s equilibrium is unaffected.
If $s_g$ is sufficiently low that $\tilde{\gamma} >
\tilde{\gamma}^*(s_g)$, the trap does not exist and $g_2$
converges to the integration equilibrium regardless of initial
economic position.
\end{proposition}

\begin{proposition}[Adverse selection through passing]
\label{prop:passing}
When individuals can reduce their personal marker salience $s_{g,i}$
at cost $\kappa(c_i)$ decreasing in class (higher-class individuals
find it easier to ``pass''), selective exit from the visible
out-group worsens the remaining group's prior:
\[
  \frac{\partial \bar{c}_g^{\mathrm{visible}}}
  {\partial (\mathrm{passing\ rate})} < 0.
\]
The individuals most able to pass---those with highest $c_i$ and
lowest natural $s_{g,i}$---are positively selected. Their departure
lowers the mean class of the visible group, which in turn raises
the discrimination rate $F(\bar{c}_g)$ against those who remain.
Individual compliance does not merely fail to improve the group
prior (the rigidity condition); it actively worsens it through
composition. The mimic man's escape harms those he leaves behind.
\end{proposition}

The formal nesting of the model within Loury's framework is developed
in Appendix~\ref{app:loury}.

% =============================================================================
\section{The Platform Extension}
\label{sec:platform}
% =============================================================================

\subsection{Client experience utility}

For a $g_1$ client $j$, experience utility at a venue with attendee
set $\A$ is:
\begin{equation}
  V_j(\A) = \alpha \cdot q + (1 - \alpha) \cdot S_j(\A)
  \label{eq:experience}
\end{equation}
where $q$ is objective quality, $\alpha \in (0,1)$ is the weight on
objective vs.\ social quality, and:
\begin{equation}
  S_j(\A) = \frac{1}{|\A|} \sum_{i \in \A,\, i \neq j}
  \E_j[c_i \mid m_i]
  \label{eq:social-quality}
\end{equation}

$S_j(\A)$ is a function of \emph{beliefs}, not facts.

\subsection{The atmosphere vector as belief architecture}

The platform produces atmosphere vector
$\mathbf{a} \in \mathcal{A}$ (pricing, location, d\'{e}cor, dress code,
marketing language, music, staff manner) serving two functions.

\textbf{Function 1: Client self-selection.}
$\Pr(\text{visit} \mid c_j, \mathbf{a}) = h(c_j, \mathbf{a})$,
increasing in the match between $c_j$ and the class $\mathbf{a}$
signals.

\textbf{Function 2: Contextual modulation of Jeffrey updating.}
$P'(g_k \mid m_i, \mathbf{a}) \neq P'(g_k \mid m_i)$. A strong
high-class atmosphere acts as a contextual prior, resolving marker
ambiguity against the stigmatised group.

\subsection{The platform's optimisation problem}

\begin{equation}
  \max_{\mathbf{a},\, \sigma} \quad
  \Pi(\mathbf{a}, \sigma) = \sum_{j \in \mathcal{D}(\mathbf{a})}
  \sigma(m_j) \cdot W_j(\mathbf{a}, \sigma) - C(\mathbf{a})
  \label{eq:platform-opt}
\end{equation}
where $\mathcal{D}(\mathbf{a})$ is the self-selecting set,
$\sigma: \M \to [0,1]$ is the admission policy, $W_j$ is willingness
to pay, and $C(\mathbf{a})$ is cost.

\textbf{Legal constraint:} $\sigma$ cannot be an explicit function of
group membership $g$. But $\sigma$ may depend on markers $m$, which
are correlated with $g$. Result: \emph{informationally mediated,
legally invisible discrimination}.

\subsection{Shadow information market}

\begin{proposition}[Shadow market intensity]
\label{prop:shadow}
The shadow information value $\Omega_j$ satisfies:
\begin{enumerate}[noitemsep]
  \item $\Omega_j$ is decreasing in $\bar{c}_g$ for $\bar{c}_g < 1/2$:
  groups with worse economic histories generate higher shadow market
  intensity.
  \item $\partial \Omega_j / \partial (1-\alpha) > 0$: clients who
  weight social quality more invest more in shadow information.
  \item $\partial \Omega_j / \partial \lambda > 0$ where $\lambda$
  indexes legal strictness: tighter anti-discrimination law raises
  the return to shadow information. The mechanism: the law constrains
  formal screening but leaves shadow cues unconstrained, degrading
  the uninformed fallback while preserving the informed outcome.
  \item $\partial \Omega_j / \partial s_g > 0$: higher marker
  salience raises shadow market intensity, because visible markers
  make group attribution easier and therefore the informational
  advantage of knowing the ``right'' venues more valuable.
\end{enumerate}
\end{proposition}

\subsection{Platform exclusion and tokenism}

\begin{proposition}[Platform exclusion of enemy group]
\label{prop:exclusion}
At the enemy group's entry prior $\bar{c}_{g_2}^0$, the platform
optimally excludes $g_2$ entirely:
$\sigma^*(m) = 0$ for all $m \in \M_{g_2}$.
\end{proposition}

\begin{proposition}[Partial inclusion equilibrium]
\label{prop:tokenism}
The platform exclusion threshold $\bar{c}_g^{\mathrm{excl}}$ satisfies
$\bar{c}_g^{\mathrm{excl}} < \bar{c}_g^{\mathrm{mid}}$: the platform begins
admitting $g_2$ members \emph{before} they have escaped the
high-discrimination basin, creating a partial inclusion zone
$[\bar{c}_g^{\mathrm{excl}}, \bar{c}_g^{\mathrm{mid}}]$ in which $g_2$ has
platform access but not social integration.
\end{proposition}

\begin{remark}
This equilibrium corresponds to what is colloquially described as
tokenism: the platform admits members of the stigmatised group before
social integration is achieved, producing visible inclusion without
the underlying belief revision that integration would require. The
result is a structural property of the Jeffrey mechanism and the
platform's profit incentives, not a claim about platform intent.
\end{remark}

\begin{proposition}[Host destabilisation]
\label{prop:spillover}
If $g_1$'s equilibrium prior $\bar{c}_{g_1}^*$ is close to
$\bar{c}_g^{\mathrm{mid}}$, the platform's admission of $g_2$ members
at $\bar{c}_g^{\mathrm{excl}}$ can push $g_1$ below its own tipping
point, triggering a class-closure episode within the host group. The
critical $g_2$ fraction $\rho_{\mathrm{crit}}$ is decreasing in the
spillover sensitivity parameter $\varepsilon$.
\end{proposition}

\subsection{Welfare analysis}

\textbf{Case 1: Beliefs accurate.} Welfare loss arises from
exclusion externality not internalised by platform or clients.
Policy: redistribution or subsidised access.

\textbf{Case 2: Beliefs inaccurate but self-confirming.}
The total welfare loss decomposes into four separable components:
\begin{equation}
  \mathcal{L} =
  \underbrace{\mathcal{L}_{\text{alloc}}}_{\text{blocked transactions}} +
  \underbrace{\mathcal{L}_{\text{info}}}_{\text{rigid priors}} +
  \underbrace{\mathcal{L}_{\text{dyn}}}_{\text{mobility suppressed}} +
  \underbrace{\mathcal{L}_{\text{platform}}}_{\text{shadow market cost}}
  \label{eq:welfare}
\end{equation}

$\mathcal{L}_{\text{alloc}}$ is the standard deadweight loss from
blocking mutually beneficial transactions between compatible agents.
It falls on both the excluded agent and the excluding agent who
foregoes a compatible interaction. This component is present in
\citet{Phelps1972}, \citet{Arrow1973}, and \citet{Loury2002}.

$\mathcal{L}_{\text{info}}$ is the novel contribution of this
framework. The rigidity condition means individual counter-examples
fail to update $F_g$; the belief system never corrects
itself even when the evidence is available. This loss falls on the
whole society because stale priors persist in all interactions, not
just those involving the excluded group.

$\mathcal{L}_{\text{dyn}}$ is the intergenerational compounding
of exclusion. Suppressed mobility lowers the group's prior across
generations (the $\delta\phi^t$ term in equation~\ref{eq:dynamic}),
which sustains exclusion, which compounds the loss. This component
grows with the distance below the critical mobility threshold
$\gamma^*$ and falls primarily on future generations of the
stigmatised group.

$\mathcal{L}_{\text{platform}}$ is the cost of the shadow
information market that arises when legal instruments constrain
formal screening. Agents invest resources in acquiring informal
information---word of mouth, coded reviews, social network
signals---that would be unnecessary in a non-discriminatory
equilibrium. This cost is borne by all agents who wish to find
their preferred social composition.

\textbf{Irreducibility result.} Three of the four components
respond to policy. $\mathcal{L}_{\text{alloc}}$ responds to
anti-discrimination law. $\mathcal{L}_{\text{dyn}}$ responds to
mobility programmes. $\mathcal{L}_{\text{platform}}$ responds to
platform regulation. The fourth component is addressed in the
following corollary.

\begin{corollary}[Irreducibility of informational welfare loss]
\label{cor:irreducibility}
Let $\Phi$ denote the class of policies that leave the soft-signal
structure of social interaction invariant: policies that operate by
changing $\gamma$ (mobility), $\sigma$ (admission), or $\delta$
(exclusion intensity), but not by converting soft markers into hard
signals. Then for all $\phi \in \Phi$:
\[
  \mathcal{L}_{\mathrm{info}}(\phi) \geq \mathcal{L}_{\mathrm{info}}^* > 0
\]
where $\mathcal{L}_{\mathrm{info}}^*$ is determined by the Jeffrey
rigidity parameter $\rho > \rho^*$ and the initial group
distribution $F_g$.
\end{corollary}

\begin{proof}[Proof sketch]
$\mathcal{L}_{\text{info}}$ arises from the gap between the
Jeffrey posterior and the Bayesian posterior, which is
the distributional gap $\rho \cdot \mathrm{dKL}(F_g \| F_g^{\mathrm{Bayes}})$---the component of belief error that
the rigidity condition prevents from updating. Policies in $\Phi$
alter $\gamma$, $\sigma$, or $\delta$ but not $\rho$ (the weight
on the group prior in the revision rule) or the group-level
conditional probabilities $F_g$. The gap is therefore
invariant to all $\phi \in \Phi$, and $\mathcal{L}_{\text{info}} >
0$ whenever $\rho > \rho^*$ (Appendix~\ref{app:loury}). \qed
\end{proof}

This result characterises the limits of formal instruments within
the class $\Phi$. It does not preclude that social movements, norm
change, or long-run cultural shifts alter the signal structure
itself---moving markers from soft to hard, or reducing the
correlation between markers and group priors. Such changes would
operate outside $\Phi$ and could in principle reduce
$\mathcal{L}_{\text{info}}$. The corollary is a theorem about
instruments, not a verdict on the world.

\begin{table}[t]
\centering
\renewcommand{\arraystretch}{1.4}
\small
\begin{tabular}{lcccc}
\toprule
\textbf{Instrument} &
$\mathcal{L}_{\text{alloc}}$ &
$\mathcal{L}_{\text{info}}$ &
$\mathcal{L}_{\text{dyn}}$ &
$\mathcal{L}_{\text{platform}}$ \\
\midrule
Anti-discrimination law   & $\downarrow$ & --- & --- & $\uparrow$ \\
Forced integration        & $\downarrow$ (temp.) & --- & --- & --- \\
Mobility programme ($\gamma > \gamma^*$) & $\downarrow\downarrow$ & --- & $\downarrow\downarrow$ & --- \\
Platform regulation       & --- & --- & --- & $\downarrow$ \\
\midrule
Irreducible               & & $\bullet$ & & \\
\bottomrule
\end{tabular}
\caption{Policy instruments mapped against welfare loss components.
$\downarrow$ = reduces; $\uparrow$ = amplifies; --- = no effect;
$\bullet$ = structurally irreducible. Anti-discrimination law
reduces allocative loss but amplifies platform loss (shadow market
effect, Proposition~\ref{prop:shadow}). Only the mobility programme
simultaneously eliminates allocative and dynamic loss above
$\gamma^*$.}
\label{tab:welfare}
\end{table}

\textbf{Case 3: Beliefs accurate and self-reinforcing.}
Policy instruments are complements, not substitutes. Forced
integration achieves equilibrium selection (shifts a cohort above
$\bar{c}_g^{\text{mid}}$) but does not alter the trap's structure.
Mobility programmes raising $\gamma > \gamma^*$ eliminate the trap
permanently through the saddle-node bifurcation
(Proposition~\ref{prop:bifurcation}). Platform regulation reduces
$\mathcal{L}_{\text{platform}}$ but has no effect on the
informational loss floor.

\subsection{Immigration, attachment tier, and closure sustainability}
\label{subsec:migration}

The platform framework extends naturally to immigration, where the
incoming foreigner's assimilation trajectory is determined by the
economic tier at which initial social attachment occurs. An immigrant
whose soft markers are immediately legible to the host society as
high-class---educational credentials in the host language,
professional dress codes, consumption patterns associated with
high-$\bar{c}_g$ groups---is assigned a high initial prior by the
host population's Jeffrey revision. Such an immigrant accesses
high-class platforms from the outset, acquires host-group high-class
markers through the endogenous marker investment mechanism
(Appendix~\ref{app:veblen}), and exits the closure trap within a
generation. An immigrant whose markers are not legible---different
cultural codes, non-transferable credentials, consumption patterns
associated with low-$\bar{c}_g$ groups in the host's prior
structure---enters the closure trap from below, regardless of actual
economic class at origin.

\begin{proposition}[Marker legibility and tier assignment]
\label{prop:legibility}
Let $m_i^{\text{foreign}}$ denote a foreign immigrant's marker
bundle. The immigrant's initial prior assignment $\bar{c}_g^{\text{init}}$
is determined by the host population's Jeffrey revision
$P'(g_k \mid m_i^{\text{foreign}})$, which depends on the legibility
of $m_i^{\text{foreign}}$ within the host marker space $\M$. If
$m_i^{\text{foreign}}$ is illegible (maps to no established group
in $\G$), the observer defaults to a low-$\bar{c}_g$ prior. Actual
economic class at origin does not enter the observer's revision.
\end{proposition}

This generates two distinct channels through which the high-income
host group's preferences over immigration composition operate.

The \emph{closure-sustainability channel}: high-income hosts have
an instrumental motive to accept immigrants pre-sorted into the
high-income tier. Such immigrants expand the high-income social
network---increasing the value of high-class platform access---while
leaving the mobility barrier for low-income groups (domestic and
foreign) intact. The group boundary shifts to absorb new arrivals
at the appropriate tier without disrupting the stratification that
sustains the closure equilibrium. Immigration of pre-sorted
high-class foreigners is closure-preserving by construction.

The \emph{competition channel}: high-income immigrants who attach
to the high-income tier compete with existing high-income hosts in
labour markets and social spaces. Whether the high-income host group
welcomes or resists same-tier arrivals depends on complementarity
versus substitutability of the immigrant's economic contribution---a
channel that is analytically independent of the closure-sustainability
motive and may cut against it.

\begin{remark}[Points-based immigration systems]
Bureaucratic immigration systems that assign points for educational
credentials, language proficiency, and professional qualifications
are formally attempts to certify high-class marker legibility in
advance of social contact. The points system converts what would
otherwise be soft, illegible foreign markers into a hard
administrative signal, enabling the host population's Jeffrey
revision to assign a high initial $\bar{c}_g^{\text{init}}$ without
direct social interaction. The empirical skilled-migration literature
documents that such systems produce faster assimilation
trajectories---consistent with the model's prediction that initial
tier assignment, not actual origin class, determines the
assimilation path.
\end{remark}

The interaction between the two channels predicts a specific
political economy of immigration preferences: high-income hosts
favour immigrants whose markers are pre-legible as high-class
(closure-sustainability dominates), are ambivalent toward immigrants
with high actual class but illegible markers (legibility friction),
and are indifferent to low-class immigrants who enter the closure
trap without disturbing high-income equilibria. Low-income hosts
resist same-tier immigration through the competition channel but
lack platform control to enforce exclusion.

% =============================================================================
\section{Discussion and Worked Examples}
\label{sec:discussion}
% =============================================================================

The model's normative stance requires brief clarification. In the
Weberian tradition of value-free social science, the analysis
identifies structural constraints on formal policy instruments
without prescribing which normative commitments should guide
political actors---the choice between competing ultimate values is
not a question social science can resolve. More fundamentally, every
agent in the model---discriminating and discriminated-against
alike---acts from within a historically constituted belief structure
that the model describes but does not rank; the Hegelian point that
reason is always socially embedded, never abstract and universal,
applies to the investigator as much as to the agents. What the model
does claim---in the tradition of Popperian piecemeal analysis---is
that well-intentioned formal interventions generate structural
unintended consequences whose identification is a precondition for
any effective political programme, whatever its normative starting
point.

\subsection{The restaurant: platform-mediated belief architecture}

An elite restaurant imposes no official racial bar and is formally
open to all who can pay. But its clients associate eliteness with
distance from lower economic backgrounds. Experience utility depends
on room composition through Jeffrey updating: clients read markers of
other diners and form posteriors over their economic class. The
restaurant, to maximise perceived quality, manages room composition
through informationally mediated instruments: dress codes calibrated
by class, reservation systems filtering by social network, location
creating transport barriers, staff trained to make certain guests
uncomfortable. Racial correlation is a consequence, not an input.
The absence of a legal channel pushes discrimination into a shadow
market of social information: coded language in reviews, word of
mouth in networks, selective social media signals among insiders.
The model predicts (Proposition~\ref{prop:shadow}) that stricter
anti-discrimination law raises the value of this shadow market,
because the law constrains formal screening while leaving informal
cues unconstrained.

\subsection{Fashion and beauty: Veblen dynamics and the foreign influence}

Fashion and beauty standards perform the function
\citet{Goffman1959} identifies as the social staging of the self: they
allow individuals to claim membership in an economic class through
visible signals. Markers that indicate health, wealth, and social
access acquire premium precisely because they encode class position.
The endogenous marker investment result (Appendix~\ref{app:veblen})
formalises this: the return to investing in class-signalling markers
is increasing in the severity of the closure trap, because the trap
raises the cost of being misclassified as low-class.

The relationship of beauty standards to foreign groups follows the
same logic and need not be driven by aesthetic preference. If a
conflict or historical episode associates a foreign group with
economic defeat or contamination---material damage to the host
society---then appearing as someone from that group becomes an
informationally costly marker to carry, because it shifts $P'(g_k)$
in observers' Jeffrey revisions toward a low-$\bar{c}_g$ group prior.
Conversely, symbols of foreign cultures that encode military or
economic dominance over the host---the Veblen trophy dynamic---raise
the value of those symbols as status markers. Whether a foreign
influence is incorporated or rejected by beauty and fashion
industries is thus ultimately a function of the economic implications
that influence carries for those who signal with it, not of
intrinsic aesthetic merit. The model provides a unified account of
both the attraction and the repulsion.

\subsection{Cultural objects and the sign-value prior: a food example}
\label{subsec:food}

The Jeffrey mechanism extends from persons to objects without
additional formal apparatus, because both occupy positions in the
same sign system grounded in economic history: the observer's
Jeffrey revision operates over sign value rather than over persons
and objects as distinct categories
\citep{Baudrillard1968,Baudrillard1972}.\footnote{This paper draws
only on Baudrillard's early work, which is compatible with a model
that grounds sign value in economic history. The later hyperreality
framework \citep[Baudrillard 1981]{Baudrillard1972}, in which signs
lose connection to any underlying referent, is not invoked and would
be in tension with the model's foundation in $\bar{c}_g$ as a sufficient
statistic of real economic history.}

A cultural object---hilsa curry, chicken feet, fermented
soybean---carries a low $\bar{c}_g$ in the host culture's prior that is
rigid to individual positive encounters by the same mechanism that
sustains person-level stereotypes. The observer's prior
$\bar{c}_g$ is reinterpreted as the probability that a
cultural object from group $g$'s cuisine is of high compatibility
with the observer's consumption preferences. The bilateral updating
structure of Section~\ref{subsec:bilateral} applies with $c_j$
denoting the object's sign position rather than a person's economic
class. The rigidity condition holds identically: a positive encounter
shifts $P'(g_k)$ for \emph{this dish} without updating
$\bar{c}_g$ for the cuisine as a category.

\begin{remark}[Object rigidity exceeds person rigidity]
The outside option $\bar{u}$ for engaging with a food is typically
higher than for engaging with a person---familiar alternatives are
readily available---which means the engagement threshold $\Delta^*$
is tighter for cultural objects than for persons. The model therefore
predicts that cultural priors should be more rigid than personal
priors without any additional assumption.
\end{remark}

\subsubsection{Heterogeneous rigidity across economic class}

Elites are more likely to cross the engagement threshold for foreign
cultural objects than the commons---not because of intrinsically
different preferences, but through two budget-constraint mechanisms.
The first is the Veblen channel: elites face diminishing
distinctiveness returns on familiar markers (their existing Veblen
signals are at risk of adoption by the class below), which
endogenously lowers their effective $\bar{u}$ for novel cultural
objects and widens $\Delta^*$. This is a consequence of
Proposition~\ref{prop:veblen} and requires no new assumption. The
second is the exposure channel, which requires one maintained
assumption.

\begin{quote}
\textbf{Assumption~E} (Exposure). Encounter frequency with foreign
cultural objects is increasing in own economic class:
$\lambda(c_i)$ is increasing in $c_i$, where $\lambda(c_i)$ is the
rate at which agent $i$ encounters cultural objects from foreign
groups. Empirical grounding: income is positively correlated with
international travel, specialist restaurant patronage, and cultural
consumption breadth.
\end{quote}

\begin{remark}[Universalised exposure does not revise the group prior]
If the budget constraint on exposure were relaxed for the entire
population---universal travel, income security, health support
abroad---exposure frequency would equalise across class. Individual
reclassification would increase uniformly (more individuals thinking
``this particular dish was surprisingly good''), but the group prior
$\bar{c}_g$ would remain rigid for the entire population.
This is the Jeffrey condition: universalising exposure universalises
individual reclassification without driving the group prior across
$\bar{c}_g^{\mathrm{mid}}$.
\end{remark}

\subsubsection{Elite adoption through the Veblen mechanism}

Once elites cross the engagement threshold and discover
compatibility, the foreign cultural object becomes a Veblen marker:
its foreignness guarantees the commons will not follow (their
$\Delta^*$ is too narrow given lower exposure and higher $\kappa$),
and the marker's value to the elite is precisely its inaccessibility
to the class below. This is fully derived from
Proposition~\ref{prop:veblen}: marker investment is individually
rational when the marker separates the investor from the class
below.

\subsubsection{Democratisation and elite abandonment}

When the barrier falls---the exotic food becomes mass-market---the
commons' $\Delta^*$ widens sufficiently for them to cross the
engagement threshold. The marker loses its distinctiveness value and
the elite abandon it, moving to the next barrier-protected cultural
object.

\begin{remark}[Congestion externality]
A full dynamic treatment requires marker value at time $t$ to
decrease in the fraction of the population that has adopted the
marker by $t$---the congestion externality formalised by
\citet{Leibenstein1950} and \citet{CorneoJeanne1997}. This mechanism
is standard in the Veblen goods literature and does not require
novel theory; its interaction with the Jeffrey architecture is
noted here without further formalisation.
\end{remark}

\subsubsection{The population-level equilibrium}

The model generates a steady-state in which a positive fraction
$\phi$ of the population privately enjoy the foreign food (their
individual $P'(g_k)$ has been updated by direct experience) while
the ambient group prior $\bar{c}_g$ for the cuisine remains
low and codes it as disgusting. The Jeffrey rigidity condition
establishes that $\bar{c}_g^t$ is invariant to $\phi^t$: the group prior
does not move regardless of how many individuals have positive
encounters.

The steady-state is characterised by two independent fixed points:
one for $\phi$ (determined by the exposure rate $\lambda$ and the
population income distribution) and one for $\bar{c}_g$ (determined by
the initial prior and the rigidity parameter $\rho$). These are
decoupled by the Jeffrey mechanism. Crucially, $\phi$ can approach 1
(everyone has privately enjoyed the food) while $\bar{c}_g$ remains at
its initial low value.

This explains a well-documented empirical pattern: foods that are
widely enjoyed by individuals from the host culture continue to be
coded as ``ethnic'' or ``exotic'' or even ``disgusting'' in ambient
cultural discourse long after individual consumption has normalised.
The model predicts this as a structural equilibrium outcome, not as
hypocrisy or inconsistency.

The Veblen mechanism introduces a second dynamic that interacts with
this population equilibrium. The elite's incentive to adopt the
marker as a Veblen signal is highest when $\phi$ is low (the marker
is distinctive) and falls as $\phi$ rises. If $\phi$ rises fast
enough through broad exposure, the elite abandon the marker before
any conceivable aggregation mechanism could shift $\bar{c}_g$. The food
cycles through three zones---disgusting, elite marker,
mass-market---without the ambient prior completing a full revision.
This is a trap that operates on cultural objects by the same
mechanism that operates on group priors for persons.

\subsection{Gender norms: inferential persistence regardless of origin}

Gender norms constitute a category of soft social markers whose
persistence the Jeffrey rigidity condition explains, regardless of
their origin or legitimacy. The model makes no claim about what
gender norms represent or whether they should be contested; it shows
only that, whatever their source---historical, institutional,
material, or cultural---the inferential mechanism that sustains class
closure sustains gender closure by the same formal route.

An individual who departs from prevailing gender norms is attributed
by observers to individual variation, not to a revision of the group
prior for their gender. This is precisely the rigidity condition of
Proposition~\ref{prop:graded}: the exception is absorbed without
updating $F_g$. Gender norms are therefore sticky even
when individual deviations are common and widely tolerated, because
the social inference process cannot propagate individual deviations
into group prior revision. The mechanism is consistent with any
account of where gender norms originate---it speaks to their
persistence, not their justification.

This observation does not advise how marginalised gender identities
should organise or resist, and it does not assess whether the
resulting equilibrium is legitimate. The positive-analytical claim
is limited to this: the Jeffrey updating rule sustains gender norms
against individual counter-examples by the same structural mechanism
that sustains racial and class stereotypes, and formal policies that
operate within the soft-signal structure of social interaction face
the same informational irreducibility as they do in those domains.

\subsection{Credential checks and protective intervention: discrimination
without adverse selection}

The model accounts for two patterns of discrimination that adverse
selection cannot explain, since neither involves repeated interaction
or self-selection dynamics.

A delivery driver in a more segregated city encounters discordant
signals: a package addressed to a culturally high-class, white-sounding
name is received by a person whose appearance suggests a low-$\bar{c}_g$
group. The driver's Jeffrey revision cannot resolve the discordance
from the available markers alone, so the driver seeks additional
markers---a credential question about the addressee's identity or
residence. This is informationally motivated marker-seeking
behaviour: the driver is attempting to sharpen $P'(g_k)$ in a
direction that resolves the conflict between the name signal and the
appearance signal. The recipient who anticipates this dynamic has an
incentive to invest in additional markers that resolve the
discordance preemptively---a particular mode of address, additional
documentation, a consumption signal at the door. This is the
endogenous marker investment of Appendix~\ref{app:veblen} in
miniature.

Restaurant staff observing a white woman seated next to a black or
Muslim man infer from the bilateral class-distance
computation under their priors that the pair is incompatible.
Their prior says these two groups do not voluntarily socialise at
close range in a high-class setting. The only interpretation
consistent with that prior is non-consensual interaction, and the
staff intervene to ``protect'' the woman. The couple sitting together
is an individual counter-example; under Jeffrey conditioning, it
shifts $P'(g_k)$ for this specific man, but it does not update
$\bar{c}_g$ for the group as a whole. The next cross-group
pair will be assessed from the same prior.

Both examples illustrate the model's general claim: discrimination
without adverse selection, driven by the rigidity of the Jeffrey
prior in the face of discordant signals. Adverse selection requires
a specific dynamic structure---repeated interaction, pool selection,
a patient-robber mechanism. Jeffrey conditioning requires only a
prior and a soft signal. The model is strictly more general.

\subsection{Compounding markers: race, gender, and institutional
credibility assessment}
\label{subsec:compounding}

The worked examples above treat a single marker dimension at a time.
The Jeffrey framework applies without modification to observers
revising over multiple marker dimensions simultaneously, and this
extension produces a qualitatively distinct result: compounding
effects that are non-additive across dimensions. The interaction
between race and gender markers in institutional credibility
assessments is a case in point, and one with direct empirical
grounding.

\subsubsection{The formal structure of compounding}

Suppose the observer is an institutional decision-maker---a police
officer, a prosecutor, a disciplinary panel---assessing an accusation
of misconduct. The observer simultaneously conducts two Jeffrey
revisions. The first is over the accused's character, using soft
markers correlated with the accused's group membership. The second
is over the accuser's credibility, using soft markers correlated
with the accuser's group membership and social position. Both
revisions are governed by the rigidity condition of
Proposition~\ref{prop:graded}: individual counter-examples shift
$P'(g_k)$ without updating $F_g$ for the group.

The compounding effect arises because the two revisions interact.
Let $\bar{c}_g^A$ denote the prior over the accused's character
(probability of innocence, say) and $\bar{c}_g^C$ denote the prior over
the accuser's credibility (probability of truthful report). The
observer's posterior assessment of the accusation's validity depends
on both simultaneously. Under Jeffrey conditioning, both priors are
held rigid against individual case-specific evidence:
$\bar{c}_g^A$ is anchored by the accused's group markers,
$\bar{c}_g^C$ is anchored by the accuser's group markers, and the joint
posterior is the product of two independently rigid revisions.

The non-additivity follows directly. If the accused carries a
race-correlated low $\bar{c}_g^A$ and the accuser carries a
gender-and-race-correlated high $\bar{c}_g^C$, the compounded
distortion of the joint posterior exceeds what either marker
dimension would produce alone. The two rigidities reinforce each
other: a downward-biased $\bar{c}_g^A$ and an upward-biased $\bar{c}_g^C$
jointly push the institutional assessment toward crediting the
accusation independently of its evidentiary content. This is not
reducible to racial discrimination alone, nor to gender bias alone;
it is the joint product of the Jeffrey architecture operating across
two marker dimensions simultaneously.

\subsubsection{The racialised harassment accusation}

The specific case that makes the non-additivity concrete is a
harassment accusation in which the accused is a Black man and the
accuser is a white woman. The observer's $\bar{c}_g^A$ is shaped by
the historical prior over Black male criminality---itself a sufficient
statistic of the group's racialised economic and legal history
(equation~\ref{eq:prior})---which is low in the host society's
inference architecture. The observer's $\bar{c}_g^C$ is shaped by the
historical construction of white female testimony as inherently
credible in this institutional context, a prior that also reflects
economic and social history rather than case-specific evidence.

Under bilateral Jeffrey conditioning, both priors are rigid to the
specific facts of the case. The accused's individual conduct, the
accuser's specific claims, and the evidentiary content of the
interaction are processed against these independently anchored
priors. The joint posterior is systematically biased in the direction
of crediting the accusation even when the evidentiary record is
ambiguous---not because any individual actor is malicious, but
because the inference architecture compounds two independently rigid
priors in a mutually reinforcing direction.

The rigidity condition explains the pattern that makes this
institutionally persistent. A case that ends in exoneration shifts
$P'(g_k)$ for this accused individual---the decision-maker may
conclude that this particular man was wrongly accused---but does
not update $\bar{c}_g^A$ for the group. The next accused Black man is
assessed from the same group prior. Similarly, a credibility
assessment that finds the accuser unreliable shifts $P'(g_k)$ for
this particular accuser but does not update $\bar{c}_g^C$ for white
female testimony as a category. The institutional prior reproduces
itself across cases by the same mechanism that sustains class and
racial stereotypes in social interaction.

\subsubsection{Symmetry and intersectionality}

The same mechanism operates symmetrically across different marker
configurations, and the direction of compounding depends on the
specific combination of group priors in play. A Black woman
reporting harassment by a white man faces a different compounding:
her race-correlated $\bar{c}_g^C$ is lower than a white woman's in the
same institutional context, while her gender-correlated $\bar{c}_g^C$
pushes in the opposite direction. The net effect on institutional
credibility assessment is not the sum of the racial and gender
components but a non-additive interaction between them---precisely
the pattern that the sociological literature on intersectionality
identifies empirically \citep{Crenshaw1991}, here given a formal
microfoundation in the Jeffrey architecture.

The model therefore does not treat gender as an isolated marker
dimension. It treats gender markers as one input into a
multi-dimensional Jeffrey revision in which race, class, and gender
all operate simultaneously and interact non-additively. The
welfare implication is that compounding discrimination is
irreducible to its component dimensions: anti-discrimination
policy targeted at the racial dimension alone, or the gender
dimension alone, leaves the interaction term untouched. Reducing
$\bar{c}_g^A$ through racial equity programmes and reducing $\bar{c}_g^C$
bias through gender equity programmes does not eliminate the
compounding distortion if each reform addresses only one dimension
of a jointly rigid inference architecture. This is a direct
extension of the informational irreducibility result
(Corollary~\ref{cor:irreducibility}) to the multi-marker setting.

\subsection{The Loury taxi driver revisited}

\citet{Loury2002} gives the example of taxi drivers who do not stop
for black passengers in high-crime areas. In his framework, this is
adverse selection: non-robbers (who are more sensitive to the cost
of waiting) exit the pool first, leaving a pool enriched in robbers,
which confirms the prior and sustains the discrimination. This is
a compelling account. In our framework it is a special case.

The Loury mechanism operates through the $\delta \phi^t$ term in
equation~(\ref{eq:dynamic}): repeated discrimination erodes the
group's economic position, which reinforces the prior, which
sustains discrimination. What the Jeffrey framework adds is the
explanation of \emph{why} the prior is rigid to individual
counter-examples. A taxi driver who does stop and finds a
law-abiding passenger updates $P'(g_k)$ for that individual---he
may think ``this person seemed respectable''---but he does not update
$\bar{c}_g$ for the group. The next passenger is assessed
from the same group prior. The rigidity condition is not an
assumption about taxi driver psychology; it is a theorem of the
updating rule that operates whenever signals are soft.

\subsection{Compliance, colonial promise, and the limits of individual
mobility}
\label{subsec:colonial}

Every historical closure system has offered a compliance
route---the promise that individuals who sufficiently adopt the host
group's markers will be reclassified and, eventually, integrated.
The \textit{\'{e}volu\'{e}} in colonial Africa was the most formally
codified instance: Africans who attained a prescribed threshold of
French or Belgian education, linguistic fluency, and cultural
assimilation were granted a quasi-civil status distinguishing them
from the colonised mass. The implicit contract was: invest in our
markers, and we will revise your group prior upward. Enough
compliance, enough time, and full integration follows.

Translated into the language of this model, the colonial promise
is a claim that the operative updating rule is Bayesian, not
Jeffrey. Under standard Bayes ($\rho = 0$, Proposition~\ref{prop:loury-nesting}),
individual compliance does update $\bar{c}_g$ for the group:
each successful \'{e}volu\'{e} is evidence that the group's class
distribution is higher than the prior implies, and the prior moves
toward the truth. Enough counter-examples, and $\bar{c}_g$ crosses the
tipping point $\bar{c}_g^{\mathrm{mid}}$; the closure trap dissolves.
This is the Loury regime: the compliance route is coherent, and
the colonial promise is structurally redeemable.

The model says the world does not operate at $\rho = 0$. Under
Jeffrey conditioning at $\rho > \rho^*$, individual marker
investment produces personal reclassification---the observer shifts
$P'(g_k)$ for this individual, attributing them to the high-class
group or to individual exceptionalism---but the rigidity condition
prevents individual compliance from aggregating into group prior
revision. The group prior $\bar{c}_g$ remains in the
high-discrimination basin regardless of how many individuals comply.
The colonial promise is structurally false in the Jeffrey regime:
not false in intention, but false as a description of how the
inference architecture actually operates. No volume of individual
compliance can drive the group prior across $\bar{c}_g^{\mathrm{mid}}$,
because the updating rule holds conditional probabilities fixed
precisely by construction.

Two features of the colonial experience that are otherwise
difficult to reconcile follow as equilibrium outcomes of the same
mechanism operating at different levels. The first is that
individual compliance often did produce genuine personal mobility:
the \'{e}volu\'{e} did achieve a different social position from
the colonised mass.
The second is that this personal mobility did not translate into
group-level liberation: the group prior moved negligibly, the
closure trap persisted, and subsequent generations of the group
began from the same discriminated baseline. The Jeffrey framework
resolves the apparent paradox without special assumptions: both
outcomes are immediate consequences of the rigidity condition. The
compliant individual is reclassified; the group is not.

The three-zone equilibrium structure of the platform analysis
(Propositions~\ref{prop:exclusion} and \ref{prop:tokenism})
formalises the position the \'{e}volu\'{e} occupied. Zone one
is full exclusion---the colonised mass, deep in the
high-discrimination basin. Zone two is compliant admission---the
\'{e}volu\'{e}, individually reclassified but collectively
disconnected from group prior revision. Zone three is full
integration---which requires the group prior itself to cross
$\bar{c}_g^{\mathrm{mid}}$ through the mobility mechanism
($\gamma > \gamma^*$), not through individual compliance. The
colonial administration, modelled as a profit-maximising platform,
admits compliant individuals into zone two because their
compliance raises the platform's atmosphere vector without
requiring genuine prior revision of the group as a whole. The
tokenism proposition is the formal counterpart of the
\'{e}volu\'{e} system: partial inclusion that falls structurally
short of integration, and that serves the platform's interests
precisely because it does.

The most diagnostically powerful historical test of the rigidity
condition is the case of colonised soldiers who fought on behalf of
the colonial power against their own people of origin. The
Tirailleurs S\'{e}n\'{e}galais, the Indian sepoys of the British
Army, the Moroccan \textit{Goumiers} in French service---compliance
in these cases was not merely cultural adoption but active military
participation in the enforcement of the colonial order. If
individual compliance could update the group prior, this extreme
form should have done so most forcefully. The historical record
shows it did not. Social admission remained partial; civic
membership was withheld; group priors in the host society were
unchanged. The rigidity condition is confirmed at the limit: even
fighting on behalf of the host identity against the identity of
origin reclassifies the individual without revising the group prior.

The model also identifies a structural cost of the compliance
route that goes beyond its failure to deliver integration.
Proposition~\ref{prop:veblen} establishes that marker investment
is increasing in trap severity: the worse the trap, the stronger
the incentive to comply. This selection mechanism extracts the
highest-investing individuals from the stigmatised group into the
compliant-admission zone, depleting the group of those most capable
of collective organisation. The compliance route is not merely
ineffective at the group level; it is actively harmful, because
it channels individual effort into personal reclassification at
the cost of collective prior revision. Colonial assimilation
policy is, on this account, a mechanism that simultaneously
produces the appearance of openness---look, the system rewards
merit and effort---while structurally preserving the closure trap
by diverting the energies that might otherwise raise $\gamma$
above $\gamma^*$.

This is not a normative verdict on those who chose the compliance
route under conditions not of their making. Individual compliance
was rational given the payoff structure of the closure equilibrium;
Proposition~\ref{prop:veblen} shows that it is the individually
optimal response to trap severity. The model's claim is only
positive: the individually rational compliance strategy is
collectively self-defeating, because the Jeffrey architecture
ensures that its aggregate effect leaves the group prior
unchanged. The route to permanent trap escape runs through
$\gamma > \gamma^*$---through structural mobility and collective
prior revision---not through individual marker investment, however
extensive.

The marker salience extension
(Corollary~\ref{cor:salience-threshold}) sharpens this analysis.
The compliance route can in principle reduce $s_g$ for markers
that are mutable: accent, dress, linguistic register, consumption
patterns. For groups with low phenotypic salience---the Irish in
England, Eastern Europeans in Western Europe---compliance
simultaneously reduces the individual's $s_g$ toward zero and the
observer's Jeffrey revision toward non-firing. In such cases, the
compliance route \emph{does} produce integration, not because it
revises the group prior, but because it renders the group marker
invisible: the individual ceases to be identified as out-group.
For groups with high phenotypic salience---racially marked
colonial subjects---compliance can modify mutable markers while
leaving the immutable phenotypic component of $s_g$ unchanged.
The Jeffrey revision continues to fire on phenotype even when
accent, dress, and cultural practice have fully converged.
The \'{e}volu\'{e}'s position is therefore doubly constrained:
the rigidity condition prevents group prior revision, and
immutable marker salience prevents individual escape from
out-group attribution. The compliance promise was structurally
false on both counts.

\subsection{Sudden legal change and the space-time mismatch}
\label{subsec:sudden-law}

The model predicts specific friction when anti-discrimination law
changes suddenly while the belief architecture remains fixed. Law
operates on the admission policy $\sigma$---a formal instrument
that is neither a spatial observation nor a temporal prior
(Remark~\ref{rem:spacetime}). It changes overnight. The Jeffrey
prior ($\bar{c}_g$, $F_g$, $\rho$) and marker salience ($s_g$)
are untouched.

Three consequences follow immediately from the existing
propositions. First, the shadow information value rises on impact:
$\partial \Omega_j / \partial \lambda > 0$
(Proposition~\ref{prop:shadow}). Discrimination moves from formal
to informal channels. Second, experience utility falls for the
host group, because social quality $S_j(\mathcal{A})$ depends on
Jeffrey-updated beliefs about new entrants, and $\bar{c}_g$ has
not moved---the host group's resistance has an
economic-inferential foundation, not just an ideological one.
Third, the result is the tokenism equilibrium of
Proposition~\ref{prop:tokenism}: visible inclusion without belief
revision.

Without simultaneous investment in $\gamma$, sudden legal change
produces an outcome that erodes political support for both
instruments: the host group experiences disutility from forced
integration that it perceives as economically costly (because the
prior has not moved), the stigmatised group experiences tokenism
rather than genuine integration, and the intensified shadow market
provides the institutional infrastructure for backlash. The
historical pattern---post-Civil Rights Act sorting in the American
South, post-Race Relations Act class-mediated exclusion in the
UK---is consistent with the model's prediction that legal change
targeting $\sigma$ without raising $\gamma$ above $\gamma^*$
redirects rather than eliminates discrimination, because the law
addresses neither the spatial observations nor the temporal priors
on which the mechanism operates.

\subsection{Digital de-anonymisation: a natural experiment on $s_g$}
\label{subsec:digital}

Anonymous digital environments---text-only forums, pseudonymous
platforms, early internet chat---are settings where $s_g \approx 0$
by design. The model predicts zero discrimination in this regime
regardless of $\bar{c}_g$: when markers are invisible, the Jeffrey
revision does not fire (Proposition~\ref{prop:salience}). The
analytically interesting transition occurs when platforms introduce
identity-revealing features: real names, profile photographs,
video. This produces a discontinuous jump in $s_g$ from near-zero
to whatever the natural marker environment supports. The model
predicts discrimination should emerge \emph{immediately} upon
marker revelation, with intensity proportional to the gap between
the revealed group's $\bar{c}_g$ and the host mean. The empirical
literature on platform discrimination is consistent with this
prediction: discrimination on rental and rideshare platforms
correlates with the introduction of photographs and real names, not
with any change in the underlying population. The mechanism is
pure spatial shock---the temporal dimension is unchanged, but the
spatial channel is switched on.

\subsection{The ``model minority'' configuration}
\label{subsec:modelminority}

East Asian immigrant communities in the US and UK occupy a region
of the $(s_g,\, \bar{c}_g)$ parameter space that the base model's
worked examples do not cover: high marker salience (phenotypically
distinct from the host group) combined with high or above-host
group mean class ($\bar{c}_{g_2} \geq \bar{c}_{g_1}$). The
expected distance gradient $\partial \E / \partial s_g = D_2 - D_1$
has a subtle observer-dependent structure in this configuration.
Since $D_2 - D_1 = (\bar{c}_{g_2} - \bar{c}_{g_1})
(\bar{c}_{g_2} + \bar{c}_{g_1} - 2c_j)$ when variances are equal,
the sign depends on the observer's class relative to the midpoint
$(\bar{c}_{g_1} + \bar{c}_{g_2})/2$. For the median host-group
observer ($c_j \approx \bar{c}_{g_1}$), $\partial \E / \partial
s_g > 0$---higher salience still increases expected incompatibility,
but the magnitude is small because $\bar{c}_{g_2}$ is close to
$\bar{c}_{g_1}$. For high-class host observers ($c_j$ above the
midpoint), the sign reverses: these observers are closer to the
out-group mean than to their own, and higher salience makes the
model minority \emph{more} attractive as an engagement partner.

The result is a split pattern: economic integration with elite
gatekeepers (who are above the midpoint and actively seek
engagement) combined with mild but persistent social friction
from the median host population (who remain below the midpoint).
This is the ``bamboo ceiling'' pattern: access to economic spaces
where high-class gatekeepers select, but exclusion from broader
social spaces where median-class members dominate. The model
predicts this is a structural equilibrium, not a transitional
state, because the spatial markers are immutable and the
observer-position structure of $D_2 - D_1$ is fixed by the
group means.

\subsection{Language revival as voluntary $s_g$ increase}
\label{subsec:language}

Language revival movements---Welsh, Hebrew, M\={a}ori---present an
apparent paradox within the model: a suppressed group deliberately
increases its own marker salience. The model predicts that higher
$s_g$ raises the mobility threshold
$\tilde{\gamma}^*(s_g)$
(Corollary~\ref{cor:salience-threshold}), so voluntary $s_g$
increase should intensify discrimination. When is this rational?
The model identifies a necessary condition: $\bar{c}_g >
\bar{c}_g^{\mathrm{mid}}$. If the group's economic position is
already above the tipping point, the closure trap does not bind
and increasing $s_g$ carries an identity value without triggering
the trap. The model predicts that successful language revival
movements should emerge \emph{after} economic convergence, not
before. Hebrew is the confirming case: revival accompanied
state-building and economic development. Indigenous language
revival in contexts of persistent economic disadvantage is the
hard case---the model predicts it is collectively costly even if
individually identity-affirming, because it raises the threshold
the group has not yet crossed. The forced marker suppression
framework (Proposition~\ref{prop:suppression}) provides the
converse: when an authority suppresses a group's linguistic
markers ($\lambda_s < s_g$), discrimination on that channel
falls but substitutes to remaining markers, and the group's
cultural identity is eroded without the group prior changing.

\subsection{Gentrification as platform optimisation}
\label{subsec:gentrification}

A gentrifying neighbourhood is a platform whose atmosphere vector
$\mathbf{a}$ is being optimised by landlords, developers, and new
businesses to attract higher-$c_j$ residents. The model's platform
extension applies directly. New amenities (coffee shops, yoga
studios, curated retail) constitute atmosphere instruments that
raise social quality $S_j(\A)$ for high-$c_j$ arrivals while
making the existing composition uncomfortable
(Proposition~\ref{prop:exclusion}). The existing residents face the
tokenism equilibrium (Proposition~\ref{prop:tokenism}): some are
retained for visible diversity while the composition shifts. The
welfare decomposition applies: amenity improvements may reduce the
allocative loss $\mathcal{L}_{\mathrm{alloc}}$, but the
informational and dynamic losses remain because $\bar{c}_g$ for the
displaced group is unchanged---they have been spatially relocated
to a lower-quality platform. When anti-displacement policies are
imposed (rent controls, affordable housing mandates), the
sudden-legal-change result (Section~\ref{subsec:sudden-law})
applies: $\sigma$ is constrained, the shadow market
(informal buyouts, service deterioration, selective maintenance)
intensifies, and discrimination redirects rather than disappears.

\subsection{Anti-colonial movements as collective action on
$\gamma$}
\label{subsec:anticolonial}

The conclusion identifies an unresolved coordination problem: how
does a dispersed and disadvantaged group coordinate on collective
mobility investment ($\gamma$) rather than individual marker
investment ($e^*$), given that the latter is individually rewarded
and the former is a public good? Anti-colonial movements are the
most important large-scale historical examples of attempted
solutions to this coordination problem, and the model can
distinguish their outcomes based on whether the movement succeeded
in raising $\gamma$ above $\gamma^*$.

\textbf{Symbiotic post-colonial relationships.} Some movements
achieved political independence---changed $\sigma$---while
maintaining the compliance channel at the cultural and economic
level. Pakistan retained British legal, military, and
administrative structures, English as the elite language, and
cricket as the national sport. Martinique remained a French
d\'{e}partement. Francophone West African states retained the CFA
franc, French military bases, and French as the official language.
In model terms, these countries changed the formal admission
instrument but continued investing collectively in $e^*$:
metropolitan marker alignment. The model predicts the tokenism
equilibrium at the national scale---formal sovereignty without
revision of the economic prior. The relationship is ``symbiotic''
precisely because the compliance channel remains open: elites from
these countries can access the metropole through individual marker
investment. The mimic man dilemma operates in full force.
Immigrants from these states face the standard model prediction:
the compliance route is available, individual reclassification is
achievable for those who invest in metropolitan markers, but the
group prior is unrevised.

\textbf{Antagonistic post-colonial relationships.} Other movements
explicitly rejected the compliance route as a national strategy.
Nehru's import substitution, promotion of Hindi, and non-alignment
constituted a deliberate attempt to build $\gamma$ independently of
the metropole while reducing dependency on metropolitan markers.
Nyerere's ujamaa was the most explicit: reject metropolitan markers
entirely, invest collectively in economic self-sufficiency.
Algeria's FLN ideology was anti-assimilation by design. The model
identifies this as the correct strategy for escaping the trap:
invest in $\gamma$, not $e^*$. But it comes at a short-term cost
---the metropole withdraws the returns to compliance, and the
transition is painful precisely because the group is attempting to
cross $\gamma^*$ without the individual safety valve of the
compliance route. Immigrants from these countries who nonetheless
pursue the individual compliance route occupy a structurally
unusual position: they are pursuing the strategy their country's
national project rejected.

\textbf{Perpetual anti-colonial narrative.} Iran after 1979
rejected the compliance route comprehensively, but $\gamma$ did not
cross $\gamma^*$ before the relationship was severed. The model
predicts that the anti-colonial narrative should persist precisely
as long as $\bar{c}_g < \bar{c}_g^{\mathrm{mid}}$: the narrative
functions as a coordination device preventing individually rational
defection back into compliance. If individuals defected into the
compliance route (cultural westernisation, emigration), an adverse
selection dynamic would operate: the most capable individuals
exit, the remaining group's $\bar{c}_g$ falls, and the trap
deepens. The perpetual
narrative makes compliance ideologically costly, keeping resources
directed toward $\gamma$ rather than $e^*$. The model predicts this
narrative will dissolve endogenously if and when $\bar{c}_g$
crosses the tipping point, at which point the trap no longer binds
and compliance becomes individually harmless.

A further prediction concerns the host country's treatment of
immigrants from each type. The group prior $\bar{c}_g$ as perceived
by the host country's observers is not purely a function of
economic history---it can be contaminated by geopolitical
antagonism. Knowledge that an individual is ``from Iran'' activates
a compound prior encoding both economic position and political
hostility. This is a case where a temporal property (ongoing
political conflict) feeds into the spatial marker (the national
label activates a prior that would not fire for an economically
identical individual from a politically neutral state). The model
predicts that immigrants from politically antagonistic states
should face more discrimination than immigrants from economically
equivalent but politically neutral states, and that this excess
discrimination should track political relations over
time---thawing relations should shift the perceived $\bar{c}_g$
toward economic fundamentals.

% =============================================================================
\section{Empirical Strategy}
\label{sec:empirical}
% =============================================================================

\subsection{Distinctive empirical prediction}

The model generates a prediction that is observationally
distinguishable from existing discrimination models.

\textbf{Non-linear convergence dynamics near
$\gamma^*$.} Standard linear mobility models predict monotone
convergence of group economic position toward the population mean.
The present model predicts non-linear dynamics: escape time from
the closure trap diverges as $(\gamma - \gamma^*)^{-1/2}$ near
the bifurcation threshold (Appendix~\ref{app:Tstar}). Groups whose
mobility parameter is close to $\gamma^*$ from below should display
slow, fitful convergence punctuated by reversals; groups well above
$\gamma^*$ should converge quickly. This power-law pattern is
observationally distinguishable from the exponential convergence
predicted by linear models.

\subsection{Data sources}

Intergenerational mobility and social network composition data: the
British Cohort Study (BCS70) and \textit{Understanding Society}
(UKHLS) provide multigenerational tracking with detailed
neighbourhood composition data for the UK. The Panel Study of
Income Dynamics (PSID) provides equivalent coverage for the US.
Cross-national comparison of group-level convergence rates against
the $(\gamma - \gamma^*)^{-1/2}$ scaling law is feasible using
harmonised OECD mobility data.

\subsection{Identification challenges}

Two endogeneity problems require attention.

\textit{Closure feedback}: discrimination in period $t$ affects
group economic position in period $t+1$, which affects the prior,
which affects discrimination. Instrument: historical economic shocks
that affected group composition before the sample period
(\citealt{NunnWantchekon2011} for Africa; colonial-era occupational
restrictions in the UK context).

\textit{Prior measurement}: $\bar{c}_g$ is latent. Use
multiple proxy measures (stated beliefs about immigrant economic
contribution, neighbourhood satisfaction, social network diversity)
and exploit variation in social desirability pressure across survey
modes as an IV for the measurement error.

% =============================================================================
\section{Conclusion}
\label{sec:conclusion}
% =============================================================================

This paper develops a formal model of social discrimination grounded
in Jeffrey conditioning---the probability-theoretic framework for
updating beliefs from soft evidence. The model generates several
results absent from the existing literature.

The observation that economic class determines social clustering is
a premise of the model, not a conclusion. What the model proves is
stronger: that the Jeffrey updating rule, applied to soft social
markers correlated with economic history, generates a
\emph{self-sustaining closure trap} as a stable equilibrium outcome,
even in the absence of racial heterogeneity. The trap is
path-dependent, intergenerationally persistent, and resistant to
individual counter-examples by construction of the updating rule.
Racial discrimination enters as a perturbation of this already-operating
equilibrium --- a second group whose mean $\bar{c}_{g_2}$
starts in the high-discrimination basin --- not as a separate
phenomenon requiring separate explanation. The closure trap requires
a minimum level of Jeffrey rigidity $\rho > \rho^*$; below this
threshold, the model collapses to a Loury-type single-equilibrium
framework where individual counter-examples erode stereotypes.

There exists a threshold level of intergenerational mobility
$\gamma^*$ below which the closure trap persists permanently and
above which it is eliminated through a saddle-node bifurcation.
Near the bifurcation, the number of generations required to escape
diverges as $(\gamma - \gamma^*)^{-1/2}$---the universal scaling of
saddle-node ghosts. This provides a formal explanation for why
apparent progress can be painfully slow even when the direction is
correct, and why programmes calibrated below $\gamma^*$ produce
measurable improvement followed by reversion.

The welfare loss decomposes into four components: allocative,
informational, dynamic, and platform. Three respond to policy.
The informational loss $\mathcal{L}_{\text{info}}$ is structurally
irreducible within the class of policies that preserve the
soft-signal character of social interaction
(Corollary~\ref{cor:irreducibility}): it arises from the updating
rule itself, not from any parameter that formal instruments can
reach. This characterises the limits of instruments, not a verdict
on the world: social movements, norm change, or long-run cultural
shifts that alter the signal structure itself could reduce this
floor.

Profit-maximising platforms create tokenism as a structural
outcome---admitting members of the stigmatised group before social
integration is achieved---and stricter anti-discrimination law raises
the return to informal discrimination networks. These are not
failures of the law but structural properties of the Jeffrey
mechanism: closing one information channel raises the return to
another.

The model identifies two qualitatively different policy instruments.
Temporary forced integration shifts a cohort above the tipping point
but does not change the trap's structure. Mobility programmes raising
$\gamma$ above $\gamma^*$ eliminate the trap permanently. The two
are complements, not substitutes. Anti-discrimination law, applied
alone, reduces allocative loss but amplifies the shadow market,
leaving the informational and dynamic losses untouched.

The marker salience extension shows that racial discrimination
cannot be reduced to economic history alone. Two groups with
identical economic histories face different closure dynamics if
their marker salience differs: a group with mutable markers (low
$s_g$) may integrate rapidly while a group with immutable phenotypic
markers (high $s_g$) remains trapped indefinitely under the same
mobility parameters. The critical mobility threshold
$\tilde{\gamma}^*$ is strictly increasing in $s_g$
(Corollary~\ref{cor:salience-threshold}), so groups with high
phenotypic salience require strictly more intergenerational
mobility to escape the closure trap. This resolves the puzzle of
why phenotypic proximity to the host group dominates severity of
historical conflict in determining integration speed, and
strengthens the compliance analysis: for high-$s_g$ groups, the
compliance route fails not only because the group prior is rigid
but because the individual cannot escape out-group attribution
regardless of how fully they adopt host-group markers.

The model identifies one further structural feature that merits
explicit statement. The individually rational response to the
closure trap---marker investment, compliance, personal
reclassification---is collectively self-defeating, because the
Jeffrey architecture ensures that individual exceptions do not
aggregate into group prior revision. The route to permanent trap
escape runs through $\gamma > \gamma^*$, not through $e^*$. The
two compete for the same scarce resources, and the Jeffrey
mechanism ensures that resources invested in individual compliance
are wasted from the group's perspective even though they are
individually rewarded. How a dispersed and disadvantaged group
coordinates on collective mobility investment in the presence of
this private incentive to defect into compliance is a coordination
problem the model identifies but does not resolve---and one that
connects naturally to the collective action literature as a
direction for future work.

% =============================================================================
% APPENDICES
% =============================================================================
\appendix

\section{The Loury Nesting}
\label{app:loury}

We introduce $\rho \in [0,1]$, the \emph{rigidity parameter}, into
the dynamic system. At $\rho = 0$, individual counter-examples fully
update the group prior $F_g$---standard Bayesian updating.
At $\rho = 1$, the group prior is completely rigid---full Jeffrey
conditioning. Combined with marker salience $s_g$, the amended
dynamic system is:
\begin{equation}
  G_{\rho,s}(\bar{c}_g, \phi) = \mu + \bar{c}_g(1 - \mu - \rho\, s_g\, \delta \phi)
  + \tilde{\gamma}\,\bar{c}_g(1-\bar{c}_g)
  \label{eq:loury-rho}
\end{equation}
The discrimination channel requires both $\rho > 0$ (Jeffrey
rigidity) and $s_g > 0$ (visible markers). Either $\rho = 0$ or
$s_g = 0$ eliminates the discrimination channel entirely, producing
$G = \mu + \bar{c}_g(1-\mu) + \tilde{\gamma}\bar{c}_g(1-\bar{c}_g)$
with a unique stable fixed point at $\bar{c}_g = 1$.

\begin{proposition}[Loury nesting]
\label{prop:loury-nesting}
There exists $\rho^* \in (0,1)$ such that:
\begin{enumerate}[noitemsep]
  \item For $\rho < \rho^*$ (Loury regime): unique stable fixed
  point at $\bar{c}_g = 1$. Individual counter-examples erode
  stereotypes. $\mathcal{L}_{\mathrm{info}} = 0$.
  \item For $\rho > \rho^*$ (Jeffrey regime): two stable fixed
  points. The closure trap exists. $\mathcal{L}_{\mathrm{info}} > 0$.
\end{enumerate}
Loury's self-confirming equilibrium is the $\rho \to 0$ limit of
this framework.
\end{proposition}

\begin{proof}[Proof sketch]
At $\rho = 0$, $\Psi_\rho > 0$ on all of $(0,1)$: no closure trap.
At $\rho = 1$, $\Psi_\rho$ dips negative: closure trap exists. By
continuity in $\rho$, there exists $\rho^* \in (0,1)$ where the
transition occurs (IVT). Under base parameters, $\rho^* \approx 0.31$.
\qed
\end{proof}

\section{Endogenous Marker Investment (Proposition~F)}
\label{app:veblen}

Agents choose marker investment $e_i \geq 0$ at cost
$c(e_i) = \frac{1}{2}\kappa_e e_i^2$. Investment shifts markers
toward high-class signals, effectively raising the $\bar{c}_g$ perceived
by others. The marginal benefit of investment is:
\begin{equation}
  \mathrm{MB}(e) = -(2(\bar{c}_g - c_j)) \cdot \theta \cdot
  (v - \bar{u}) \cdot (1 + \xi \phi^*)
  \label{eq:MB}
\end{equation}
where $\theta > 0$ is marker sensitivity and $\xi > 0$ captures the
amplification from trap severity $\phi^* = F(\bar{c}_g^L)$.

\begin{proposition}[Endogenous marker investment]
\label{prop:veblen}
The equilibrium marker investment is:
\begin{equation}
  e^* = \frac{(2(c_j - \bar{c}_g)) \cdot \theta(v-\bar{u})(1+\xi\phi^*)}
  {\kappa_e}
  \label{eq:estar}
\end{equation}
with the following properties:
\begin{enumerate}[noitemsep]
  \item $e^* > 0$ iff $c_j > \bar{c}_g$: only agents above
  the class midpoint invest in markers.
  \item $\partial e^*/\partial \phi^* > 0$: more severe closure trap
  generates more marker investment.
  \item The non-negativity constraint $e^* \geq 0$ binds for agents
  with $c_j < \bar{c}_g$, producing a separating equilibrium:
  high-class agents invest in markers, low-class agents pool at
  zero.
\end{enumerate}
\end{proposition}

This is the formal model of \citet{Veblen1899}'s conspicuous
consumption: marker investment performs the function of signalling
class distance from the stigmatised group. It is also
\citet{Goffman1959}'s impression management formalised as an
optimisation problem within the Jeffrey framework.

\section{Intergenerational Persistence ($T^*$)}
\label{app:Tstar}

\begin{corollary}[Critical slowing near the bifurcation]
\label{cor:Tstar}
Near $\gamma = \gamma^*$, the number of generations required to
escape the closure trap from initial position $\bar{c}_g^0$ scales as:
\begin{equation}
  T^*(\gamma) \sim C \cdot (\gamma - \gamma^*)^{-1/2}
  \quad \text{as } \gamma \to \gamma^{*+}
  \label{eq:Tstar}
\end{equation}
where $C = \pi / \sqrt{|a \cdot b|}$, with $a = \partial \Psi /
\partial \gamma$ and $b = \frac{1}{2}\partial^2 \Psi / \partial
\bar{c}_g^2$ evaluated at the bifurcation point.
\end{corollary}

The $-1/2$ exponent is the universal scaling of saddle-node ghost
passage times. It follows from the normal form of the bifurcation:
$\Psi \approx a(\gamma - \gamma^*) + b(\bar{c}_g - \bar{c}_g^*)^2$ near the
collision point. The map spends $O((\gamma - \gamma^*)^{-1/2})$
iterations in the ghost region where $|\Psi|$ is $O(\gamma - \gamma^*)$.

\textbf{Empirical prediction.} A programme raising $\gamma$ by $0.01$
above $\gamma^*$ requires approximately 65 generations to achieve
integration---multiple human lifetimes. The distribution of
convergence times across groups should follow a power law with
exponent $-1/2$ near the bifurcation. ESS repeated cross-sections
could test this at the country level.

% =============================================================================
% BIBLIOGRAPHY
% =============================================================================

\begin{thebibliography}{99}

\bibitem[Akerlof(1970)]{Akerlof1970}
Akerlof, G.~A. (1970).
\newblock The market for ``lemons'': Quality uncertainty and the market
mechanism.
\newblock \textit{Quarterly Journal of Economics}, 84(3):488--500.

\bibitem[Akerlof and Kranton(2000)]{AkerlofKranton2000}
Akerlof, G.~A. and R.~E. Kranton (2000).
\newblock Economics and identity.
\newblock \textit{Quarterly Journal of Economics}, 115(3):715--753.

\bibitem[Arrow(1973)]{Arrow1973}
Arrow, K.~J. (1973).
\newblock The theory of discrimination.
\newblock In O.~Ashenfelter and A.~Rees (eds.), \textit{Discrimination in
Labor Markets}. Princeton University Press.

\bibitem[Baudrillard(1968)]{Baudrillard1968}
Baudrillard, J. (1968).
\newblock \textit{Le Syst\`{e}me des objets}.
\newblock Gallimard, Paris.
(Translated as \textit{The System of Objects}, Verso, 1996.)

\bibitem[Baudrillard(1972)]{Baudrillard1972}
Baudrillard, J. (1972).
\newblock \textit{Pour une critique de l'\'{e}conomie politique du signe}.
\newblock Gallimard, Paris.
(Translated as \textit{For a Critique of the Political Economy of the Sign},
Telos Press, 1981.)

\bibitem[Becker(1957)]{Becker1957}
Becker, G.~S. (1957).
\newblock \textit{The Economics of Discrimination}.
\newblock University of Chicago Press.

\bibitem[Bernheim(1994)]{Bernheim1994}
Bernheim, B.~D. (1994).
\newblock A theory of conformity.
\newblock \textit{Journal of Political Economy}, 102(5):841--877.

\bibitem[Bonilla-Silva(1997)]{BonillaSilva1997}
Bonilla-Silva, E. (1997).
\newblock Rethinking racism: Toward a structural interpretation.
\newblock \textit{American Sociological Review}, 62(3):465--480.

\bibitem[Bordalo et al.(2016)]{BordaloEtal2016}
Bordalo, P., K.~Coffman, N.~Gennaioli, and A.~Shleifer (2016).
\newblock Stereotypes.
\newblock \textit{Quarterly Journal of Economics}, 131(4):1753--1794.

\bibitem[Coate and Loury(1993)]{CoateLoury1993}
Coate, S. and G.~C. Loury (1993).
\newblock Will affirmative-action policies eliminate negative stereotypes?
\newblock \textit{American Economic Review}, 83(5):1220--1240.

\bibitem[Coetzee(1999)]{Coetzee1999}
Coetzee, J.~M. (1999).
\newblock \textit{Disgrace}.
\newblock Secker \& Warburg.

\bibitem[Currarini et al.(2009)]{CurrariniEtal2009}
Currarini, S., M.~O. Jackson, and P.~Pin (2009).
\newblock An economic model of friendship: Homophily, minorities, and
segregation.
\newblock \textit{Econometrica}, 77(4):1003--1045.

\bibitem[Crenshaw(1991)]{Crenshaw1991}
Crenshaw, K. (1991).
\newblock Mapping the margins: Intersectionality, identity politics, and
violence against women of color.
\newblock \textit{Stanford Law Review}, 43(6):1241--1299.

\bibitem[Corneo and Jeanne(1997)]{CorneoJeanne1997}
Corneo, G. and O.~Jeanne (1997).
\newblock Conspicuous consumption, snobbism and conformism.
\newblock \textit{Journal of Public Economics}, 66(1):55--71.

\bibitem[De~Finetti(1964)]{DeFinetti1964}
De~Finetti, B. (1964).
\newblock Foresight: its logical laws, its subjective sources.
\newblock In H.~E. Kyburg and H.~E. Smokler (eds.),
\textit{Studies in Subjective Probability}. Wiley, New York.
\newblock (Original: La pr\'{e}vision: ses lois logiques, ses sources
subjectives, \textit{Annales de l'Institut Henri Poincar\'{e}}, 1937.)

\bibitem[Fanon(1952)]{Fanon1952}
Fanon, F. (1952).
\newblock \textit{Peau noire, masques blancs} [Black Skin, White Masks].
\newblock \'{E}ditions du Seuil, Paris. English translation: Grove Press, 1967.

\bibitem[Galbraith(1952)]{Galbraith1952}
Galbraith, J.~K. (1952).
\newblock \textit{American Capitalism: The Concept of Countervailing Power}.
\newblock Houghton Mifflin, Boston.

\bibitem[Galbraith(1958)]{Galbraith1958}
Galbraith, J.~K. (1958).
\newblock \textit{The Affluent Society}.
\newblock Houghton Mifflin, Boston.

\bibitem[Galbraith(1967)]{Galbraith1967}
Galbraith, J.~K. (1967).
\newblock \textit{The New Industrial State}.
\newblock Houghton Mifflin, Boston.

\bibitem[Goffman(1959)]{Goffman1959}
Goffman, E. (1959).
\newblock \textit{The Presentation of Self in Everyday Life}.
\newblock Doubleday.

\bibitem[Jeffrey(1965)]{Jeffrey1965}
Jeffrey, R.~C. (1965).
\newblock \textit{The Logic of Decision}.
\newblock McGraw-Hill.

\bibitem[Jeffrey(1983)]{Jeffrey1983}
Jeffrey, R.~C. (1983).
\newblock Bayesianism with a human face.
\newblock In J.~Earman (ed.), \textit{Testing Scientific Theories}.
University of Minnesota Press.

\bibitem[Kendall(2002)]{Kendall2002}
Kendall, D. (2002).
\newblock \textit{The Power of Good Deeds: Privileged Women and the Social
Reproduction of the Upper Class}.
\newblock Rowman \& Littlefield.

\bibitem[Leibenstein(1950)]{Leibenstein1950}
Leibenstein, H. (1950).
\newblock Bandwagon, snob, and Veblen effects in the theory of consumers'
demand.
\newblock \textit{Quarterly Journal of Economics}, 64(2):183--207.

\bibitem[Loury(2002)]{Loury2002}
Loury, G.~C. (2002).
\newblock \textit{The Anatomy of Racial Inequality}.
\newblock Harvard University Press.

\bibitem[Naipaul(1967)]{Naipaul1967}
Naipaul, V.~S. (1967).
\newblock \textit{The Mimic Men}.
\newblock Andr\'{e} Deutsch.

\bibitem[Nunn and Wantchekon(2011)]{NunnWantchekon2011}
Nunn, N. and L.~Wantchekon (2011).
\newblock The slave trade and the origins of mistrust in Africa.
\newblock \textit{American Economic Review}, 101(7):3221--3252.

\bibitem[Phelps(1972)]{Phelps1972}
Phelps, E.~S. (1972).
\newblock The statistical theory of racism and sexism.
\newblock \textit{American Economic Review}, 62(4):659--661.

\bibitem[Spence(1973)]{Spence1973}
Spence, M. (1973).
\newblock Job market signalling.
\newblock \textit{Quarterly Journal of Economics}, 87(3):355--374.

\bibitem[Veblen(1899)]{Veblen1899}
Veblen, T. (1899).
\newblock \textit{The Theory of the Leisure Class}.
\newblock Macmillan.

\bibitem[Du~Bois(1935)]{DuBois1935}
Du~Bois, W.~E.~B. (1935).
\newblock \textit{Black Reconstruction in America}.
\newblock Harcourt, Brace and Company, New York.

\bibitem[Walley(1991)]{Walley1991}
Walley, P. (1991).
\newblock \textit{Statistical Reasoning with Imprecise Probabilities}.
\newblock Chapman and Hall, London.

\bibitem[Weber(1922)]{Weber1922}
Weber, M. (1922).
\newblock \textit{Wirtschaft und Gesellschaft}.
\newblock Mohr, T\"{u}bingen.
(Translated as \textit{Economy and Society}, University of California
Press, 1978.)

\end{thebibliography}

\end{document}
